%%%%%%%%%%%%%%%%%%%%%%%%%%%%%%%%%%%%%%%%%%%%%%%%%
\section{The WSU System Architecture}
%%%%%%%%%%%%%%%%%%%%%%%%%%%%%%%%%%%%%%%%%%%%%%%%%

This section presents the WSU system architecture. Similar to the structure of \cite{CSA-IA}, the WSU system architecture is presented in three different layers, i.e. the processes (Section~3.1), supported by the hardware (Section~3.2) and software (Section~3.3) infrastructure. The processes section is ``agnostic'', in the sense that it does not make any reference to specific SW or HW components. It simply describes the (future) complete ALMA operational chain. On the other hand, the HW and SW infrastructure sections do name the various components of the ALMA operational chain. The WSU ALMA system architecture can be visualized in Fig.~\ref{fig:proc_hw_sw_wsu}.

%%%%%%%%%%%%%%%%%%%%%
\begin{figure}[htb]
\centering
\includegraphics[width=0.9\textwidth]{Figures/Processes_HW_SW-WSU.png}
\caption{Overview of the WSU ALMA system architecture. An additional process, related to the commissioning of new observing modes, is not present in this figure as it happens in parallel with the WSU processes.
%, rather than being part of the ALMA operational chain.
\label{fig:proc_hw_sw_wsu}} 
\end{figure}
%%%%%%%%%%%%%%%%%%%%%

\clearpage


%%%%%%%%%%%%%%%%%%%%%%%%%%%%%%%%%%%%%%%%%%%%%%%%%
\subsection{The WSU Processes}
\label{sec:wsu_processes}

%Thodori to add text about this section:  

This section covers the complete ALMA operational chain, starting with the proposal submission and project preparation,  followed by the scheduling and data acquisition, corrections for atmospheric and instrumental effects, data processing and quality assurance, and, finally, ingestion of the calibrated data and images into the ALMA Science Archive. The flow of these processes is depicted in the middle panel of Fig.~\ref{fig:proc_hw_sw_wsu}.

As a guide to the rest of this chapter, the general hardware and software infrastructure is indicated in the left and right panels of Fig.~\ref{fig:proc_hw_sw_wsu}. 
As these processes are agnostic to the type or name of the software or hardware used, there is no specific mention to what is known in the ALMA jargon as subsystems (i.e. the software components that support the processes).   

As previously mentioned, a description of the current ALMA processes is provided in Section~3.1 of~\cite{CSA-IA}. This section only describes how the current processes will evolve, to cope with the data obtained in the WSU context. 

%%%%%%%%%%%%%%%%%%%%%%%%%%%%%%%%
\subsubsection{Proposals Submission and Project Preparation}
%%%%%%%%%%%%%%%%%%%%%%%%%%%%%%%%
During the proposal preparation phase, the observatory provides resources to the community to prepare and submit observation programs. All details about the proposal and project preparation phase can be found in Section~3.1.1 of~\cite{CSA-IA}.  

In the WSU era, similar project preparation activities to those currently in place are anticipated. Nevertheless, the SW subsystem used by the observatory for (a) proposal tracking and management (Section~3.3.1 of~\cite{CSA-IA}), and (b) generating the proposal statistics produced in each Cycle, will need to adopt to the WSU-required modifications of the ALMA Project Data Model (APDM)\index{ALMA Project Data Model} structure~\cite{apdm}. 

The creation of new states or substates and updating transitions between states (project, Observing Unit Set (OUS), or Scheduling Blocks (SB)\index{Scheduling Block(s)}), might be required, as they will be necessary for the Group OUS processing implementation \cite{dp}. Duplication and resubmission policies should be compatible with the WSU capability (see Section~4.1.1.2 of~\cite{CSA-IA}). Automation will be needed in calibration implementation during Phase~2 generation and the technical assessment.

Additionally, the process of the array configuration selection might need to be re-evaluated, as the beam-size will vary significantly over the wide range of the observed frequency.


%The main relevant tools during proposal preparation are documentation, Observation Tool (OT) for designing and submitting projects, Phase 1 Manager (Ph1M) for observatory tracking and managing proposals, DPS for simulation, help desk, ALMA archives, duplicate checking (metadata and python scripts) and other complementary tools and resources. Duplication and resubmission policies are compatible with WSU capability, and the checking tool also follows it.

%Projects are submitted through the OT, which will be upgraded for the WSU (see Section~\ref{subsec:OT}). After the CfP is closed, Phase 1 SBs are automatically generated that contains basic information on all Science Goals for proposals. With this information, submitted proposals are checked against the archive and current projects in the queue to identify duplicates and resubmissions. During this step, a PDF of each proposal is also generated by the OT and distributed by the PHT to the reviewers. Large Programs (LPs) are evaluated for technical and schedule feasibility.

%Once the review process is completed, PHT delivers the scientific rankings to the Queue Building (QB) process, that delivers the assigned grades for each project for approval by the Director’s Council. Upon grades approval, the notification letters are sent to the PIs, and Phase 2 SB preparation process begins. The identified duplicates are descoped from the queue, and the identified resubmitted Science Goals (SG) are confirmed with PIs to avoid duplicated observations. After automatic Phase 2 SB generation with the OT, the Phase 2 Group reviews and modify as required the SBs, including modifications following a technical assessment. Calibration activities are triggered when required and preparation process ends when SB status are set to ready for observations. These Phase 2 generation processes are also expected to be automated in the WSU era.


%An observing project is composed of a hierarchy of structures, with the project at the top level, followed by OUSs, and SBs at the lowest level (Fig.~\ref{fig:projectStructure}). All the possible states from initiation to completion, and the transitions between them, are defined as a life cycle.


%%%%%%%%%%%%%%%%%%%%%
%\begin{figure}[t]
%\centering
%\includegraphics[width=0.85\textwidth]%{Figures/projectStructure.png}
%\caption{Block diagram of an Observing Project from the point of view of the observation preparation.
%\label{fig:projectStructure}} 
%\end{figure}
%%%%%%%%%%%%%%%%%%%%%


%%%%%%%%%%%%%%%%%%%%%%%%%%%%%%%%
\subsubsection{Scheduling and Data Acquisition}
%%%%%%%%%%%%%%%%%%%%%%%%%%%%%%%%

The scheduling and data acquisition process manages the execution of SBs on each array for Science Operations. While each subsystem and software will need to cope with the WSU, the basic process itself is as described in Section 3.1.2 of~\cite{CSA-IA} and no major changes are anticipated. However, the scheduling and the data acquisition process need to accommodate to the parallel deployment and the WSU in terms of the resources of arrays such as Array Elements (AEs), ATAC\index{Advanced Technology ALMA Correlator}/Total Power GPU Spectrometer (TPGS)\index{Total Power GPU Spectrometer}, bands, available time for the Science Operations, and the data rate. Similar to project preparation, the algorithm for selecting appropriate SBs for a given array configuration during scheduling, based on the angular resolution required by the PI, may need to be updated to cope with the WSU. Adjustments to the scheduling may also be needed to prioritize completion of GOUSs. 

Once SBs are to be executed, the main difference in the WSU era will be the need for Control to handle a mix of legacy and upgraded receiver cartridges within a single band across the array, and to manage communication with both new and old devices. See Section~\ref{handling_legacy_upgraded_devices} for more details.

% In the scheduling process, the list of SBs is provided by the Dynamic Scheduler Algorithm (DSA) to AoD with priority based on inputs such as scientific ranking score, APDM database, Shiftlog database, AQUA QA0 database, Life cycle database, phase-tracker and PWV (weather conditions), available hardware resources in the data-acquisition infrastructure (ATAC and TPGS, operational AEs and bands, and photonic references), and operational policies and priorities. 

% Based on the outcome of DSA, the most suitable SB for the current conditions is selected for execution and is sent for execution via the Science Software Requirements (SSR) subsystem that acts as interface between Scheduler and Control. Each observation, i.e. execution of an SB, generates an Execution Block (EB), whose are stored in the ALMA Archive in ASDM format (see Section~\ref{subsec:ASDM-APDM}). An SB is observed multiple times until the aggregated EBs meet the
% original sensitivity requirement, currently evaluated on execution time on source corrected by real weather and
% the number of AEs in the actual observations. Once this condition is fulfilled, the MOUS containing the SB is
% considered as ``Fully-Observed'' and proceeds to data processing.
% The data acquisition process includes the following: 
% \begin{itemize}[itemsep=-0.5em]
%     \item Resource allocation by creation of arrays.
%     \item Assessment of weather condition and available bands.
%     \item Selection of appropriate SBs in accordance with scheduling outcome.
%     \item SB queuing and management.
%     \item Monitoring array status and performances.
%     \item Logging of all observing activities and recording problems during
% the operations.
% \end{itemize}


%%%%%%%%%%%%%%%%%%%%%%%%%%%%%%%%
\subsubsection{Calibrations}    % Todd
%%%%%%%%%%%%%%%%%%%%%%%%%%%%%%%%
\label{sec:calibrations}
\paragraph{Antenna Calibration}
\newcommand\offlineCalibration{3.1.3.3}

In the WSU, the antenna calibration tasks will continue to be needed 
to maintain the readiness and efficiency of all array elements for PI science 
observations (see Section~3.1.3.1 of~\cite{CSA-IA}).  The fundamental nature of ALMA as
a reconfigurable interferometer
means that all the normal antenna integration processes will be executed each time an antenna is relocated, including pointing model verification, delay adjustment and position determination.  The performance requirements in these areas will remain the same, but per-EB (or per-session calibration for polarization observations) will require a faster cadence for antenna position corrections and/or modifications to the current QA0 processes.
The components of the instrumental delay model will be adapted to the new IF arrangement with only two baseband pairs.  

\paragraph{Online Calibration}

The concept of ``online'' calibration refers to the calculation and application of quantization corrections and the normalization of visibility spectra, along with the calculation of atmospheric calibration tables that can be used in offline processing to convert the stored visibilities into antenna temperature-calibrated spectra (see Section~3.1.3.3 of~\cite{CSA-IA}).   
% beginning of review comment 13
% The following commented text was replaced by the subsequent text.
%All of these techniques will be implemented in the WSU system, but there will be some changes in the details.  
These techniques have been implemented in the current system, but changes will be needed for the WSU.
% end of review comment 13
Quantization corrections will
be applied for all Imaging Correlation (IC) modes, in contrast to the current system in which the BLC does not apply them in FDM spectral setups.  Visibility normalization will continue to be performed at full spectral resolution, but the atmospheric calibration tables computed online might not be at full resolution due to 
processing resource limitations, in which case the full solution will be computed offline instead.

\paragraph{Visibility Data Calibration}

%%%%%%%%%%%%%%%%%%%%%
\begin{itemize}
\item {\bf WVR correction:}
WVR corrections will continue to be essential to achieve ALMA's science goals in the WSU, and the method and cadence of their calculation and application will generally remain as it is now -- offline (in the pipeline) for imaging correlation modes, and online (in the correlator, using the WVR calibration coefficients supplied by telescope calibration SW) for VLBI/Beamforming mode (see 
Section~\offlineCalibration\/ of~\cite{CSA-IA}).  The implementation details of the latter will be adapted to the new correlator architecture.
%%%%%%%%%
\item{\bf Bandpass Calibration:} 
The basic concept of bandpass calibration using an observation of a strong point source
to solve for the antenna-based spectral response in amplitude and phase will
remain the same in the WSU (see Section~\offlineCalibration\/ of~\cite{CSA-IA}).  With fewer analog components in the IF signal
chain, the temporal stability of the WSU instrumental bandpass may
be better than today, but it will need to be measured and demonstrated 
during commissioning that the current cadence (once per EB) will still be sufficient.  
One major change envisioned is that the spectral resolution of the bandpass scan may be limited to 1-2\,MHz if the science target spectral resolution is $<$2\,MHz, in order to
%One major change envisioned is that the spectral resolution of the bandpass scan may be limited to  be $\gtrsim 2$~MHz 
in order to reduce data volume and processing time (particularly for QA0, see section 3.4.4 of \cite{da}) while not adversely affecting science data products in terms of SNR achieved on the data cubes~\cite{bandpassMemo}.  With this 
limitation, combined with the improved digital efficiency of the WSU, the current moderate length scans ($5-30$~min, depending on array and frequency band) will be sufficient to achieve the critical SNR of 50 in many cases, with the automated channel pre-averaging invoked  by the pipeline solver in other cases.  But due to the difference in spectral resolution, and hence spectral setup, the pipeline will need to be adjusted to apply a bandpass measured from one spw to a different spw, matching by the baseband name and spectral window number contained in the spw name.  For broad spws at low frequency, the spectral index across the bandwidth may vary significantly enough to require a more involved treatment in the pipeline.
%%%%%%%%%
\item{\bf Gain Calibration:}
The methods of measuring temporal changes in the antenna-based complex gain (amplitude and phase), via in-band or band-to-band calibration, will continue in the WSU (see Section~\offlineCalibration\/ of~\cite{CSA-IA}).  The improvements in sensitivity, due to the general increase in 
aggregate correlated bandwidth, will likely be traded for selecting gain
calibrators closer to the science target, and thus achieving better overall
removal of atmospheric phase errors while also providing more less sensitivity
to antenna position errors.
%There might be a period of time that the bandwidth switching style of in-band calibration becomes important for projects seeking high spectral resolution in an era of data rate caps.
%%%%%%%%%
\item{\bf Flux Density Calibration:}
Solar system objects, in particular Uranus \cite{c10_thb,Farren2021}, will remain the fundamental flux density standard in the WSU, in that they anchor the flux density scale of the grid calibrators
%that are routinely observed in multiple bands to track their variations in flux density and spectral index with time 
(see Section~\offlineCalibration\/ of~\cite{CSA-IA}). The flux service, which interpolates these measurements in time and frequency, will continue to be an essential component of ALMA in the WSU, with further improvements in the heuristics foreseen, particularly in the new low frequency bands.  
Flux calibration shall be performed such that the resulting flux calibration is linear with respect to
the flux scaling, and be able to re-scale the flux for individual EBs prior to imaging as necessary.
This will allow the per-session calibration to be run as close as possible to the time the data
were acquired, while the fluxes can be re-scaled before imaging, if necessary, either because the
source catalog values changed or to equalize the flux between EBs.
Spws with wider fractional bandwidth enabled by the WSU will need improved treatment in the pipeline to  maintain the 
accuracy of the flux scale. 
%%%%%%%%%
\item{\bf Polarization Calibration:}
Measurement of the antenna-based polarization calibration of the WSU analog signal path will, as now, require observing a polarized quasar over sufficient parallactic angle rotation to solve for the necessary quantities (see Section~\offlineCalibration\/ of~\cite{CSA-IA}).  To achieve sufficient angular coverage, multiple EBs typically need to be observed, and in a session block that prevents changes to the IF attenuator settings in order to avoid changing the relative path delay and phase of the two linear feeds.  The WSU use of wide
fractional bandwidths in the new low frequency bands may require solving
%the DPS to be able to solve 
for the frequency dependence of the calibrator’s Stokes Q and U when computing the XY-phase estimate.
\end{itemize}
%%%%%%%%%%%%%%%%%%%%

\paragraph{Total Power (TP) Data Calibration}
WSU TP observations on-source will be performed in the current
on-the-fly (OTF) mode, during which the antenna(s) scan the sky at a constant speed, while the TPGS continuously captures spectra with a short integration duration.
Amplitude calibration will be achieved by observing a nearby, emission-free
position, to compute a differential signal in temperature units then scaling by the gain factor (Jy/K) of the antenna.   Solar observations will continue to use 
the broadband continuum detectors, which will accept the wider bandwidth of
the expanded WSU IF range (see Section~\ref{sec:IFConditioner}) compared to the 2~GHz baseband filters that precede the
SQLDs in the current system.


%%%%%%%%%%%%%%%%%%%%%%%%%%%%%%%%
\subsubsection{Quality Assurance Level 0/1 (QA0/QA1)}
%%%%%%%%%%%%%%%%%%%%%%%%%%%%%%%%
\label{sec:QA01}
% QA0 and QA2 are the main quality assurance processes routinely executed on ALMA datasets. QA0 is done on every EB and QA2 on every MOUS. QA2 implies calibrating and imaging the ALMA data. It is triggered only after all the SB executions needed to reach PI's goals have been collected. Currently, imaging is done only per MOUS. Details regarding QA0 can be found in Section~3.1.4 of~\cite{CSA-IA} and Section~11.2 of~\cite{c10_thb}.

% In the ALMA Observing Tool,  PI's science goals are transformed into  observational units called Group OUSs (see Section~3.1.1 of~\cite{CSA-IA}). To reach a science goal, observations in different Arrays (12-m, 7-m, TP) or configurations may be needed. Consequently, a GOUS is split into MOUSs, each of them corresponding to a different array or array configuration. Every MOUSs contains one SB. Several executions of the same SB may be needed to achieve the required sensitivity, meaning that a MOUS may contain several EBs. 

The first quality check along the data chain is QA0, which is performed on each EB as soon as it is observed.  It is intended to identify data that are are clearly of inadequate quality to merit further processing and assessment, and to refine the execution fractions used by the scheduling subsystem. In order for the quality of EBs to be accurately assessed, they need to be fully calibrated by running them through all the evaluation heuristics of the official processing SW, in contrast to the less comprehensive method currently in place. This improvement in the workflow is desirable independent of the WSU.
Therefore, QA0 in the WSU era is expected to include running the full calibration pipeline of all EBs immediately following their observation, but constrained to a per-session basis for polarization projects.\footnote{Currently, polarization observations require multiple scans of the polarization calibrator through parallactic angle and are performed through a series of EBs observed with exactly the same instrumental setting, termed a session. Session continuity is broken by any reset of an antenna or correlator.} 

An initial assessment of the EBs will still be necessary, to filter out the most obvious ones not complying with required calibration content or quality standards. Yet, the QA0 status will be finalized only after full EB calibration. %, with the calculation of its execution fraction (EF): the estimation of how well the EB matches its expectation (as defined by the ALMA OT from PI's goals) in terms of sensitivity. The correct estimation of the EF is fundamental to evaluate whether further observing of the same SB is needed or not and this is why it constitutes  the final step of QA0.


QA1 was intended to be the process by which the observatory monitors the behavior of the instrument over time, looking for trends and triggering fixes as needed. It should be noted that to date the project has never defined or put into place any formal QA1 workflows. If QA1 is deemed essential for WSU, the necessary workflows will need to be clearly defined and put into production.

%%%%%%%%%%%%%%%%%%%%%%%%%%%%%%%%
\subsubsection{Data Processing and the Quality Assurance Level 2/3 (QA2/QA3) Process}
%%%%%%%%%%%%%%%%%%%%%%%%%%%%%%%%
The calibrated data produced in QA0 will be used to proceed to the imaging step on which the QA2 evaluation will be performed. QA2 is the second tier, and normally the last tier, of quality assurance that each PI dataset goes through. Its aim is to ensure that all applicable quality standards are met, and that the PI science requirements are satisfied. With WSU, the first step in QA2 will commence immediately following QA0 by automatically pipeline calibrating each EB (or session for polarization observations). The results of this calibration will be reviewed, if QA scores indicate this is needed; and the calibration tables (or calibrated visibilities themselves) stored for use later in the QA2 process.

%The aim of QA2 is to ensure that the data products meet the PI requirements. Such evaluation is expected to occur through a QA2 process that uses quantitative parameters calculated by the ALMA Pipeline.

%, so that when the calibration of an EB ends, we will both finalize its QA0 and have a QA2 assessment of its calibration.

\textbf{With the WSU, a major change will be that interferometric imaging will be performed per Group Observing Unit Set (GOUS), rather than per MOUS} (\cite{dp}). This means that it will be executed on the combination of all the MOUSs in the necessary configurations and arrays, to satisfy a science goal (see Section~3.1.1 of~\cite{CSA-IA}). 
In general, full science quality images\footnote{We note that today not all PIs receive full science quality images at the MOUS level: 19\% of 12-m data sets in Cycle~7 were either produced with less than optimal qualities or not generated at all \cite{memo623}.} of all targets and all spws at the MOUS level are not expected to be delivered to PIs for computational, storage, and efficiency reasons. A subset of representative products are likely to be generated in order to support MOUS level QA2 evaluation, and these could be delivered to the PI; see also the special cases listed in Section~3.6.6 of \cite{dp}.
%For computational, storage and efficiency reasons, delivering images produced at MOUS level is not generally foreseen. GOUS imaging will be triggered when all MOUS in a group meet certain QA criteria~\cite{dp}. 
The key advantage of this imaging workflow is that 
the GOUS product will be more directly matched to achieving the user's science goals, while the ongoing operational cost should be comparable to the current operational cost of assessing and delivering the individual MOUS products as ALMA currently does \cite{dp}. However, group level QA2 criteria need to be defined, and their relation to observation (MOUS) level criteria considered.
 
Several checks will be put in place in the quality assurance phases prior to imaging, i.e. at QA0 level.  Once all the executions in a MOUS have been collected and calibrated, the MOUS as a whole will be evaluated to check whether it will achieve its expected sensitivity and angular scales, using sample images and cubes (i.e., not the full science products). This process is needed to assess any cases that would fail QA at the MOUS level quickly enough to schedule more observations when necessary. In particular, the evaluation must be fast enough to do so even for the highest data rate projects. 
 

%With such an approach, the current distinction between QA0 and QA2 will become less strict, as EB calibration occurs at the interface of these two processes,  being at the same time the last step of QA0 and the first one of QA2. 

%Quality Assurance level 2 (QA2) deals with evaluating whether PI's goals have been met. The first elements on which QA2 will be performed are the calibration products coming out of the QA0 process per session. This will be a major change in QA2 in the WSU era, since currently calibration occurs only after a MOUS is fully observed. 

%When a PI needs to sample quite different angular scales and possibly recover the zero-scale flux, their science goal is split into different MOUSs, each corresponding to a different array or array configuration, grouped in one GOUS. In the current operational model, imaging is performed per MOUS, hence per array and per configuration. With the WSU, PIs will receive GOUS imaging products, which are  the actual ones matching their requests. 

A last step of QA2 will need to include a final check that the images/cubes indeed achieve the PI's goals. As much as possible, all such checks will need to be automated, based on quantitative parameters. 

%A schematic representation of the workflow is presented in Figure~\ref{fig:dpoverview}.

%\begin{figure}[p]
%\includegraphics[width=\linewidth]{Figures/DP_Figure1b.png}
%\caption{An overview of the WSU data processing workflow. Stated %timescales do not apply to the highest volume datasets.} 
%\label{fig:dpoverview}
%\end{figure}


%To summarize, the WSU QA2 will be split into
%\begin{itemize} 
%\item[1)] a calibration process that we may call QA-cal, triggered as the final step of QA0, 
%\item[2)] an intermediate evaluation of whether goals were achieved in terms of sensitivity and angular-scale sampling (QA-obs) performed on calibration products at MOUS level,  and 
%\item[3)] an imaging process (QA-img) that will happen at GOUS level.
%\end{itemize}

The operational model for the QA3 process, which aims at correcting issues found in the data delivered to the PIs and can be triggered by users or by the Observatory itself, is described in \cite{dp}.

%%%%%%%%%%%%%%%%%%%%%%%%%%%%%%%%
\subsubsection{Archive Ingestion}
%%%%%%%%%%%%%%%%%%%%%%%%%%%%%%%%
\label{archiveIngestion}

Two major changes for Archive ingestion are including GOUS-level science products and per-EB calibration. The former is allowed for in the current design. Support for per-EB calibration will need to be added. Finally, any possible ALMA Science Data Model (ASDM)\index{ALMA Science Data Model} changes (see Section~\ref{ASDMchanges}) will impact archive ingestion.

%%%%%%%%%%%%%%%%%%%%%%%%%%%%%%%%
%\subsubsection{Commissioning of Observing Modes}
\subsubsection{Commissioning of Pipeline for New Observing Modes and Improved Heuristics}
%%%%%%%%%%%%%%%%%%%%%%%%%%%%%%%%

From the outset, the basic functional work of PLWG members and developers means that manual interaction with data is explicitly required to facilitate building the processing schemes for new observing modes into the ALMA DPS Pipeline tasks or to build new ALMA DPS Pipeline tasks themselves, and for the continual development/modification of the ALMA DPS Pipeline tasks. There must be functionality within the ALMA DPS to allow the manual interaction and manipulation of visibility data at the lowest level - i.e. extracting visibilities into e.g. Python arrays. It therefore goes without saying that the ALMA DPS code, at some level -- e.g. inside Pipeline specific areas/tasks -- will be open access for modification and/or could allow plug-in additions, such as heuristic, QA, or plotting codes at appropriate points (within a task/stage function). A helper framework\footnote{E.g. high-performance code would be under the hood, but elements, such as heuristic use for statistics, would be within the framework of tools at a top level, e.g. a standard-deviation calculation to process any parameterized sub-set of data, field, scan, or spw, or a standard plot function routine that can be passed a variety of data.} functionality inside the ALMA DPS core code could be envisaged to facilitate the correct formulation of said heuristics, plotting, or QA code development. Continued discussion with NRAO's Data Management \& Software Division (DMS) will be required\footnote{More details about DMS can be found in Section 4 of~\cite{dp}}, as the application layer (i.e. lower-level API on which the high-level pipeline-like tasks are built) is designed, to ensure that it is as accessible as possible to the wide range of ALMA internal users (commissioning scientists, developers, PLWG scientists, DRMs).

However, many full datasets in the WSU era will not be manageable for commissioning and development following the current testing and acceptance scenarios and a modified workflow/working process needs to be devised, which ultimately means the ALMA DPS must be designed in such a way as to enable that workflow. Considering the worldwide geographical location of PLWG members and their various compute resources, a case can be built for easy and fast data transfer, shared/bespoke compute resources, and transferable Pipeline runs/jobs (with associated data) that can be run repeatedly per-stage indefinitely without reliance of a priori stages. To explore these further:
%%%%%%%%%%%%%
\begin{itemize}
    \item It would be preferential that the ALMA DPS tasks are built to be able to process a ``minimal function block" of data, e.g. bandpass calibration and all heuristics, plots, QA etc. are confined to running on only the bandpass field, or any stage can simply run on the required field, scan, and selected spw without the rest of the whole observational data. 
    \item PLWG members and developers may not have the compute resources to iteratively test any modifications, even when considering the above point. The Pipeline team (developers and PLWG) will need access to specific dedicated compute facilities and data storage (via network access).  In addition, the process of instantiating an unreleased SW branch on that remote hardware, gathering the appropriate data and SW dependencies, and running and debugging on that remote hardware needs to be as automated/scripted and accessible as possible.
    \item ALMA DPS Pipeline stages/tasks must  be rewindable to a state where the run is still valid, i.e. before any issues/bugs occur in a subsequent stage. Any given stage should be able to be rerun indefinitely. It would be highly desirable if on-the-fly code changes could be made and the stage rerun, or at least a Pipeline run would be re-loadable after code changes and restarted at exactly the same point. In the case of working with other PLWG members and DRMs, and for problem investigations, the entire Pipeline job and state should be transportable and could therefore be restarted at another location/institute.  
\end{itemize}
%%%%%%%%%%%%%

\newpage

%%%%%%%%%%%%%%%%%%%%%%%%%%%%%%%%%%%%%%%%%%%%%%%%%
\subsection{The WSU Hardware Infrastructure}
%%%%%%%%%%%%%%%%%%%%%%%%%%%%%%%%


%%%%%%%%%%%%%%%%%%%%%%%%%%%%%%%%%%%%%%%%%%%%%%%%%
%\subsubsection{A Preliminary WSU Architecture}% with Parallel Deployment }
%%%%%%%%%%%%%%%%%%%%%%%%%%%%%%%%%%%%%%%%%%%%%%%%%

Based on the outcome of various development studies led by the Executives in the past decade, and the efforts of three initial working groups~\cite{2ndgen,fed,scwg}, established by the AMT, a new strawman ALMA system architecture fulfilling the needs of ALMA 2030 was developed. This strawman system architecture has been evolving in the last few years as the concept has become more mature and different WSU-related Development Projects with definite plans have been approved by the ALMA Board. This now preliminary system architecture continues to evolve through the steady design progress of the various Development Projects, following the Product Assurance requirements in force within ALMA~\cite{pa_req}, and their interactions via the drafting of the Interface Control Documents (ICDs). The ALMA System Technical Requirements have also been flown down to some of the lower-level products currently under development, with corresponding subsystem requirements reviews, and this has in some cases allowed the identification of additional necessary requirements at system level.

This section provides an overview of the preliminary ALMA 2030 system architecture from a hardware point of view, focusing on four key areas where major changes to the existing system architecture are expected, 1) the analog signal chain including the digitizers, 2) the near real-time digital signal processing, including the Data Transmission System (DTS), 3) the correlator and the spectrometer, and 4) the site infrastructure. 
%An analysis of system requirements and their flow-down to lower-level products is only meaningful if an underlying system architecture is present. A new ALMA system architecture fulfilling the needs of ALMA 2030 has not been established yet but based on the outcome of the various development studies by the Executives made over the past decade and the efforts of the three working groups, established by the AMT, substantial input is available to define a reasonable working assumption. This preliminary system architecture is a starting point towards a new baseline, the process for establishing this new formal baseline shall comply with the Product Assurance requirements in force within ALMA~\cite{pa_req}. Design reviews not only at system level but also at subsystem level are essential for a controlled and structured process to obtain the new ALMA 2030 system architecture.

%This section provides a first concrete proposal in defining this ALMA 2030 System Architecture and can be considered the core of the Conceptual System Design Description. It has taken the current system architecture as a basis, while changes have been made primarily following the recommendations from the working groups~\cite{2ndgen,fed} within constraints set by JAO and AMT to minimize the upgrades to areas where it is clearly justified and essential to reach the ALMA 2030 objectives and to minimize the operational impact. On this last topic, it should be emphasized that a concept of ``parallel deployment'' will be followed in order to minimize the disruption to science observations. Thus, operations using the old (``legacy'') system can be interleaved (in time) with operations using the upgraded system.  In addition, operations on the legacy digitizer/DTS/correlator system that combine some legacy system components with some upgraded system components can be expected. This aspect is mostly likely foreseen for the Front-End receivers (CCA + WCA).  It is important to note that the existing software system (both online and offline, see sections 5.7 - 5.10) will need modifications to enable and support the parallel deployment of the hardware.  A freeze on new software features for the legacy ALMA system will likely be needed in order to implement these critical WSU software updates.

%The introduction of the preliminary architecture will focus on three key areas where major changes to the existing system architecture are expected, 1) the analog signal chain including the digitizers, 2) the near real-time digital signal processing, including the Data Transmission System (DTS), and 3) the correlator and total power spectrometer.  Subsequent sections will describe the upgrade of each major component of the signal chain, including the control software and upstream/downstream software and computing systems.



%%%%%%%%%%%%%%%%%%%%%%%%%%%%%%%%%%%%%%%%%%%%%%%%%
\subsubsection{Signal Chain from the Front-Ends to the Correlator/Spectrometer}
%%%%%%%%%%%%%%%%%%%%%%%%%%%%%%%%%%%%%%%%%%%%%%%%%
Fig.~\ref{fig:analog-signal-chain} provides a simplified signal chain block diagram presenting the proposed analog signal chain. 
The general trend in modern high performance receiver designs is to digitize the analog signal as early as possible in the signal chain and perform the signal processing in the digital domain if technically and economically viable. The digitizer largely determines the boundary between the analog and digital domains.
%In modern high performance receiver designs, there is a general trend to digitize the analog signal as early as possible in the signal chain and, as much as technically and economically viable, perform signal processing in the digital domain. Where this boundary between analog and digital domains is set, is largely determined by the digitizer. 
Given the requirements of ALMA, chiefly in terms of instantaneous bandwidth and signal quality, a direct digitization of the RF signal is not yet feasible. However, with the substantial progress made in digitizer technology, sampling and quantization of the first IF signal has become within reach of ALMA~\cite{Quert_final}, even considering the stretch goal of 4$\times$ the current instantaneous bandwidth. This technology breakthrough made it feasible to simplify the analog signal chain and obviate the requirement for a second IF in the BE subsystem.

%%%%%%%%%%%%%%%%%%%%%
\begin{figure}[t]
%\hspace*{-0.3cm}
\includegraphics[width=0.90\textwidth]{Figures/DownConv_LOmix_DelayCorr.png}
%{Figures/Analog-signal-chain-Mar2024.png}  
% figure was updated to remove TFB and correlator clock, add timestamping, and change WVR processing to the more general term of pipeline processing. - Todd, 27Feb2024;  modified again to add "or Spectrometer" - Todd 22Mar2024
\caption{Simplified ALMA 2030 preliminary system block diagram focusing on the analog signal chain.  The frequencies included in this diagram are just working examples and shall not be understood as requirements.  The actual LO1 frequency range is $31-938$~GHz and the actual RF range is $35-950$~GHz, each being a function of the receiver band.
\label{fig:analog-signal-chain}} 
\end{figure}
%%%%%%%%%%%%%%%%%%%%%

\clearpage

%%%%%%%%%%%%%%%%%%%%%
\begin{figure}[t]
\includegraphics[width=\textwidth]{Figures/Antenna-to-Correlator-Diagram.png}
\caption{Simplified ALMA 2030 system block diagram focusing on the near-real time digital signal path.\label{fig:antennaToCorrelator}} 
\end{figure}
%%%%%%%%%%%%%%%%%%%%%


Fig.~\ref{fig:antennaToCorrelator} introduces the changes to the near-real time digital signal processing chain. Important changes to the current system architecture are the following:
\begin{itemize}[itemsep=-0.5em]
    \item \underline{Correlator architecture}: the current Frequency division\,/\,Cross Correlation\,/\,Fourier transform (FXF) architecture of the BLC will be replaced by a ``poly-phase'' Frequency Division\,/\,Fourier Transform\,/\,Cross correlation (FFX) architecture. This fundamental change has been extensively analyzed by correlator and system experts and there is consensus that an FFX architecture is the most effective solution to implement the ALMA 2030 requirements. This architecture choice for the WSU was confirmed in~\cite{ffx_decision}.
    \item \underline{Data Transmission System (DTS)}: the current DTS is a custom designed, real-time optical transmission system. The capacity of this DTS cannot cope with the increased data rate from the antennas, $\sim$1~Tb/s of raw data for each antenna, and a suitable solution needs to be identified. Given the costs and risks involved in developing a custom system, preference was given from the beginning to identify a Commercial-Off-The-Shelf (COTS) solution. ALMA 2030 will be in the fortunate position to benefit from the substantial technological and commercial progress made in very high-speed optical data links. The Laboratoire d’Astrophysique de Bordeaux team at University of Bordeaux and the Signal Chain Working Group have studied this subject and recommended to consider very high-speed optical Ethernet-like solutions, e.g. 100~GbE or 400~GbE, for this purpose. 400~GbE equipment has entered the market and is expected to be a very mature, commercial product by the time that the ALMA DTS upgrade will be made; thus, 400~GbE is being used in the DTS preliminary design (Section~\ref{sec:DTS}).
    
    An important consequence of adopting a COTS Ethernet solution for the DTS is that it will no longer be a real-time transmission channel since it is packet-based. This issue can be resolved in the ALMA 2030 system architecture by accurately time-stamping the data samples produced by the digitizers. By making this metadata available to the correlator subsystem, it will still be feasible to perform the correlation function. In fact, this correlation function does not need to be done in real-time, each data sample has an unambiguous time stamp, but in practice it will be done by the upgraded system architecture in near real-time.
    Accurately time stamping samples at the antennas increases the complexity of the digitizer electronics but this additional effort is well compensated by the use of a COTS DTS solution compared to a custom one.
    \item \underline{Correlator physical location}: Primarily due to the capabilities of an Ethernet-like DTS, it has now become feasible to increase the physical distance between antennas and correlator subsystem. Currently, the whole Correlator subsystem, including the Tunable Filter Banks (which for high spectral resolution modes provide a ``first F''), is co-located at the AOS Building. With the new DTS, the future Correlator subsystem can be moved to the OSF site, which is highly desirable from an operational point of view, including ease of maintenance. The new DTS makes it even feasible to distribute the location of the Correlator subsystem among different physical locations. One option that was under consideration was to locate the first F of the Correlator subsystem at the antennas, close to the digitizers. The justification for this option was that formatting and time stamping of the sampled data needs FPGA devices that have excess processing capability which can also be used to perform the first F function. A further analysis of this correlator architecture concluded that on technical and operational grounds the preferred solution was to have no first F at the antennas~\cite{IDTsubbanding} as was also recommended by various working groups.   The top two considerations were that preventative actions, maintenance procedures, and future technology upgrades are all easier to manage and accomplish when all the correlator components are located in the same room; and that the correlator software team would only need to interact with hardware delivered by a single hardware group, simplifying project and schedule management. The option to have no first F at the antennas was also endorsed by the AMT considering technical and programmatic impacts~\cite{ffx_decision}.
\end{itemize}




%With this preliminary definition of the architecture, the concept of parallel deployment of the WSU systems (see Section~\ref{sec:parallelDeployment}) alongside the existing ALMA systems is summarized in the diagram in Fig.~\ref{fig:parallelDeploymentBlockDiagram}, which illustrates the two separate paths of the signal chain.  %Science observations with the existing system could be interleaved in time with commissioning observations with the new system, for example on nighttime vs. daytime periods, with the corresponding change of software versions.



%%%%%%%%%%%%%%%%%%%%%
%  THIS FIGURE HAS BEEN MOVED TO ParallelDeploymentConcept.tex
%\begin{figure}[htb]
%\includegraphics[width=\textwidth]{Figures/ALMA2030_Concept_of_Parallel_Deployment.png}
%\caption{Parallel deployment concept enabling observations with either the existing ALMA systems or the WSU systems. Note that the proposed interleaved 40 GSps digitizers require a digitizer clock rate of 20~GHz.\label{fig:parallelDeploymentBlockDiagram}} 
%\end{figure}
%%%%%%%%%%%%%%%%%%%%%



%%%%%%%%%%%%%%%%%%%%%%%%%%%%%%%%%%%%%%%%%%%%%%%%%
\subsubsection{Front-End}
%%%%%%%%%%%%%%%%%%%%%%%%%%%%%%%%%%%%%%%%%%%%%%%%%
\label{sec:frontend}
\paragraph{Receivers}

Following the antenna mirror surfaces and cabin membrane, the ALMA signal chain begins at the FE receivers. Most of the current receivers (Bands $3-10$) are based on SIS mixers operating near 4~K with a noise performance that meets the original ALMA requirements~\cite{sys_reqs}, in many cases by a significant margin~\cite{fed}.   However, the SIS mixers, as well as the CLNAs that follow the mixers in the signal chain, are based on the technology from 10 years ago (or more). Receiver noise temperature could be improved immediately by upgrading the existing mixers (and in some cases the CLNAs) with current state-of-the-art SIS junction and semiconductor transistor devices. Additionally, for the higher Bands ($6-10$), several groups have investigated higher current density ($J_c$) superconducting junction technology and preliminary results indicate the possibility of manufacturing SIS mixers with wider IF and RF bandwidths, as well as flat noise temperature performance across the full bands~\cite{ESOfeconf}.  As a result of these ongoing developments in the field, an ALMA working group devised updated receiver performance goals (for noise temperature, sideband rejection, passband gain variation, polarization purity, and beam squint), accounting for the on-array performance of existing ALMA receivers and typical effective sky temperatures at the ALMA site. The technical goals are intended to be ambitious but attainable on the 2030 timescale and are summarised in~\cite{fed}.  The ALMA System Technical Requirements~\cite{sys_reqs} have been updated according to these goals.  The actual IF bandwidth and noise temperature requirements may need to be adjusted for each band to reflect the practical limitations of technology.

Band 2 will be the first receiver band to deliver the wider bandwidth goal of ALMA 2030.  Band 2 will be a 2SB receiver employing advanced wide-RF CLNAs and a downconverter with a $2-18$~GHz IF~\cite{b2cre}, for which six pre-production cartridges~\cite{b2} were produced and the corresponding MRR was passed in September 2023. The production project for the rest of the cartridges is ongoing.

The first receiver band to be upgraded as part of the WSU is Band 6, which began in January 2022 as ``Phase 1'' ALMA development project to build a prototype upgraded warm cartridge assembly (including lower noise YIG LOs and wider bandwidth room temperature IF amplifiers) and a prototype upgraded cold cartridge assembly (including superconducting mixers and CLNAs) using SIS mixers constructed with advanced junction fabrication technology~\cite{b6up}. The IF bandwidth per sideband and per polarization will be at least 12~GHz ($4-16$~GHz) with a goal of 16~GHz.  The PDR is expected later in 2024 while initial tests of the prototype at the ALMA site are scheduled to be in 2025.

In November 2023, the ALMA Board approved the development of an upgraded Band 8 receiver unit addressing the goals of the WSU, which would be followed up by a production proposal for the rest of the cartridges. The IF bandwidth per sideband and per polarization will be at least 14~GHz ($4-18$~GHz) with a goal of 16~GHz. Improvements in receiver noise temperature will be investigated during the preliminary design phase.

The tentative order in which the rest of the bands will be updated is discussed in the WSU implementation plan which describes three phases of the upgrade~\cite{imp_plan}.  It is envisioned that Bands 9 and 10 will ultimately be replaced by 2SB receivers.  


%%%%%%%%%%%%%%%%%%%%%%%%%%%%%%%%%%%%%%%%%%%%%%%%%
\paragraph{First Local Oscillator (LO1)}
%%%%%%%%%%%%%%%%%%%%%%%%%%%%%%%%%%%%%%%%%%%%%%%%%
The LO1 reference signal is distributed by the Central LO photonic system via fiber optic cables to each antenna. The Central LO photonic system consists of a Master Laser (ML) locked to a master frequency standard and a Slave Laser (SL) that is offset locked to the ML. The output of these two lasers is combined and transmitted to the antennas. At the antennas, the desired LO1 reference signal is generated by photomixing the two infrared laser carrier tones to produce a tunable frequency reference ($31.0-121.7$~GHz tunable, band-dependent) for use by the antenna based LO1 systems. The Central LO system also contains line length correctors to compensate for thermal expansion and mechanical stresses on the fibers. No changes are needed in this aspect of the LO1 system to support the WSU.

In the FE Warm Cartridge Assembly (WCA) for Bands $2-10$, a high-Q YIG Tuned Oscillator (YTO) forms the fundamental oscillator in a band-dependent sub-band in the $12.2-27.4$~GHz range that is multiplied up to the desired LO1 “driver” frequency that is phase locked to the incoming LO1 phase/frequency reference in an offset phase-locked loop whose other input is the First Local Oscillator Offset Generator (FLOOG) synthesizer frequency. The phase locked WCA output or LO1 “driver” frequency is multiplied up by a cryogenic LO multiplier in the cold cartridge to provide the LO signal for the SIS mixers. In Band 1, the YTO operates over a range of $31-40$~GHz and requires no multiplication. The required tuning ranges of the YTOs were designed to lie above the highest frequency of the ALMA IF range (12~GHz) for all receiver bands. As shown in Fig.~\ref{fig:YTOranges}, the broader IF range of the WSU will begin to encompass the YTO ranges. For instance, in the case of Band 6, where the YTO frequency range is $12.22-14.77$~GHz, the YTO will always fall in the broader IF range. In part for this reason, the Band 6 receiver upgrade project~\cite{b6up} has proposed a new AMC architecture which will employ a new higher frequency YTO with a frequency range double the current value, eliminating the possible direct leakage of this strong signal into the science IF range. The YTOs for the other bands can all be parked at or above 16.9~GHz, to avoid interfering with the receiver IF. However, as the IF range of upgraded receiver bands push beyond $\sim$17~GHz, similar measures should be considered in the WCAs of other Bands undergoing upgrades, in order to eliminate this source of RFI generated by the FE system.

%%%%%%%%%%%%%%%%%%%%%
\begin{figure}[htb]
\includegraphics[width=0.9\textwidth]{Figures/FE-LO-YLO-freq-ranges.png}
\caption{Plot showing the ALMA FE LO YTO frequency ranges of all receiver bands (horizontal lines) with respect to the current IF range of most 2SB bands ($4-8$~GHz, black--to--blue vertical dashed lines), and the potential WSU IF ranges of 2 or 4 to 16, 18, or 20~GHz (grey-or-black--to--red vertical dashed lines). Credit: G.~Siringo and K.~Saini.  Note that the existing YTO for Band~6 (brown) would fall within any extended IF range that includes $\sim 12-15$~GHz.  This potential problem will be alleviated by implementing the $\times 3$ AMC architecture in the Band~6v2 WCA upgrade (pink). The YTOs for all other bands could, in principle, be parked at $ >16.9$~GHz when in standby modes, and thus would not interfere with a large portion of the expanded IF range of a WSU receiver. Band~5 (purple) has the lowest maximum YIG frequency and will need special care prior to upgrading the WCA of that band.\label{fig:YTOranges}} 
\end{figure}
%%%%%%%%%%%%%%%%%%%%%

In addition, the harmonics of the YTOs can become a concern.  Currently, the Band~5 LO is parked (when not in use) at 13.33~GHz, to keep its 6th harmonic at 80~GHz, i.e. below the 84~GHz lower edge of the Band~3 IF range. But for Band~2, there is no such solution, as the 6th harmonic of the Band~5 YTO ranges from $79.98-101.50$~GHz, which means that it will always fall somewhere in the RF range of Band 2.   Because Band 5 is not scheduled for upgrade until the final phase of the WSU~\cite{imp_plan}, mitigation strategies will need to be explored, such as adjusting its YTO parking frequency depending on the spw locations in the active Band~2 science SB. 


%%%%%%%%%%%%%%%%%%%%%%%%%%%%%%%%%%%%%%%%%%%%%%%%%
\paragraph{IF Switches}
%%%%%%%%%%%%%%%%%%%%%%%%%%%%%%%%%%%%%%%%%%%%%%%%%
A relatively minor but crucial part of the WSU is the upgrade of the existing 10-way IF switch inside the FE Assembly to support the wider bandwidth IF signals.  Also, the addition of a new 2-way isolating IF power splitter, hereafter called the ``WSU IF Splitter'', will be a key hardware element to enable “parallel deployment”, which will accommodate the operational requirement of minimizing the Observatory down time during the installation and commissioning of the WSU.  In other words, parallel deployment will provide the flexibility of performing either science observations with the current system or commissioning observations with the upgraded system (at different times, not simultaneously). The bandwidth of both of these IF components shall match the stretched ALMA 2030 goal and provide at least 4 times the current IF bandwidth.  The Signal Chain Working Group has proposed a design concept for such an upgraded IF chain in~\cite{scwg}.


%%%%%%%%%%%%%%%%%%%%%%%%%%%%%%%%%%%%%%%%%%%%%%%%%
\subsubsection{Back-End}
%%%%%%%%%%%%%%%%%%%%%%%%%%%%%%%%%%%%%%%%%%%%%%%%%

\paragraph{IF Conditioner}
\label{sec:IFConditioner}
A functional block diagram for the proposed IF signal chain for the IF Conditioner, which accepts the output of the Front End IF switch, is provided in Fig.~\ref{fig:IFswitch}, including the various amplifiers, variable attenuators, filters, couplers, and continuum detectors.    To avoid the lengthy and, frequency dependent, lossy coaxial cable transporting the analog IF$_{1}$ signals from FE Assembly to BE analog rack, 
one option is to co-locate the digitizers with the upgraded IF switch and other analog components placed inside the FE Assembly.  However, the working group on ``IF Switch unit and Digitizer possible configurations and locations'' \cite{IFSwitchLocation} produced a different recommendation. 
The impact of the proposed location needs to be assessed from a cost/benefit perspective at future design reviews
%\sout{(see Section 6.2)} 
including Electromagnetic Compatibility (EMC) concerns (see also Section~\ref{DTSconstraints}).  

As an upgrade to the existing Square Law Detector (SQLD) circuit, whose response is limited by the 2~GHz baseband filters, the WSU broadband continuum detectors will cover the wider WSU IF range (i.e. the whole sideband) as set by the anti-aliasing filter preceding the power splitter ahead of the digitizer.  With one per sideband per polarization channel, these detectors will be used to determine the IF attenuator level settings, to measure broadband \Tsys, and to observe the Sun in continuum mode for Total Power (TP)\index{TP, Total Power} solar science observations.   Additional narrower band detectors might also be added to provide flexibility for engineering purposes. The digital data from these detectors will be collected by the control software into in a separate binary stream sent to the Archive.

\begin{figure}[htb]
\includegraphics[width=\textwidth]{Figures/IFSwitchConditioner.png}
\caption{Functional block diagram taken from~\cite{scwg}, for a single polarization/sideband output channel, including a new ``WSU IF Splitter'' that can connect the selected Front-End signal (drawn from 1 of 10 bands) for use by either the legacy Back-End (top path) or the new WSU Back-End (bottom path).\label{fig:IFswitch}} 
\end{figure}

%%%%%%%%%%%%%%%%%%%%%%%%%%%%%%%%%%%%%%%%%%%%%%%%%
\paragraph{Digitizers}
%%%%%%%%%%%%%%%%%%%%%%%%%%%%%%%%%%%%%%%%%%%%%%%%%
\label{digitizerSection}
As introduced in Section~\ref{sec:digitizers}, commercial digitizer technology has substantially advanced over the past decade and a device has been identified that matches the ALMA 2030 objectives very well. The preferred device  (Fig.~\ref{fig:Micram}) has been designed and is marketed by the German company Micram, acquired by Keysight in March 2022. The key performance parameters of this device are:

%%%%%%%%%%%%%%%%%%%%
\begin{figure}[htb]
\includegraphics[width=\textwidth]{Figures/MicramADC2.png}
\caption{Digitizer assembly based on Micram/Keysight 6-bit ADC module.\label{fig:Micram}} 
\end{figure}

\begin{itemize}[itemsep=-0.7em]
    \item Technology:			        \quad Infineon SiGe BiCMOS process ($f_{\rm T}$ / $f_{\rm max}$: 170 / 250~GHz)
    \item Digitizer architecture:		\quad dual interleaved flash converter
    \item analog input BW:		       \quad 20~GHz
    \item Effective sampling rate:	   \quad 40~GSps
    \item Number of bits/sample:	   \quad 6
    \item Power dissipation:		   \quad $10-13$~W
\end{itemize}
%%%%%%%%%%%%%%%

Four of these digitizers, two units (LSB and USB) per polarization, will be installed in each antenna to support in access of 16~GHz bandwidth per polarization and sideband. 
More information, including quantization efficiency and the Effective Number Of Bits (ENOB) can be found in~\cite{Quert_final}. The IF conditioner preceding the digitizer will include an anti-aliasing filter with an upper limit cutoff frequency ($-1$~dB point) as close as possible to 20~GHz to support upgraded receiver bands that are targeting the $4-20$~GHz IF range.  Prior to delivery to the DTS, the digitized data packets will be timestamped so that they can be correctly aligned by the new correlator (Section~\ref{sec:ATAC}) and the spectrometer (Section~\ref{sec:TPGS}).  Design of the packet format will need to be coordinated among all groups.

%%%%%%%%%%%%%%%%%%%%%%%%%%%%%%%%%%%%%%%%%%%%%%%%%
\paragraph{Frequency/Timing Jitter Requirements}
%%%%%%%%%%%%%%%%%%%%%%%%%%%%%%%%%%%%%%%%%%%%%%%%%
The current BE specifications document (Section~5.71 of~\cite{BEreqs}), defines the Timing Event (TE / 48~ms) requirements, including jitter, as delivered at each antenna by the Back-End subsystem. TE jitter is dependent on the transmission medium used, for Low Voltage Differential Signaling (LVDS) the specification value is $\leq$ 0.8~ns while for RS-422\footnote{ANSI/TIA/EIA-422-B Electrical Characteristics of Balanced Voltage Differential
Interface Circuits.} the specification is $\leq$ 8~ns. This current performance is not stable enough for time stamping purposes with the proposed ALMA 2030 architecture using an asynchronous DTS.
The jitter for time stamping purposes of the digitized data, assuming an effective sampling rate of 40~GSps, must be better than 2.5~ps (assuming 98\% coherence at band edges). Achieving this jitter requirement for antenna-based oscillators and timing circuits, synchronized to the Central Local Oscillator (CLO), should be feasible considering that the current LO2 phase jitter, which is phase locked to the 125~MHz reference frequency, already meets a requirement of less than 1~ps (Req. BEND-04522-00/RT~\cite{BEreqs}). The LVDS technology itself is not a limiting factor for this matter since a jitter contribution of less than 1~ps can be achieved.


%%%%%%%%%%%%%%%%%%%%%%%%%%%%%%%%%%%%%%%%%%%%%%%%%
\paragraph{Establishing and Maintaining Absolute Time}
%%%%%%%%%%%%%%%%%%%%%%%%%%%%%%%%%%%%%%%%%%%%%%%%%
At present, when the ALMA Array Real Time Machine (ARTM) of the control system starts up, it performs a one-time alignment of the TEs with the 1 Pulse-Per-Second (PPS) signal from the Global Positioning System (GPS) so that there is a timing event at the top of each minute and every 48~ms thereafter. Regarding the method of maintaining absolute time for the WSU, there are two options: 1) use the existing method, or 2) use a COTS Precision Time Protocol (PTP) server linked with the GPS.  In the existing method, the control system sends a message to clients (such as the Antenna Bus Master (ABM)\footnote{ABM is the real-time computer that controls all devices in an antenna.}): ``the time at the next TE is HH:MM:SS.sss'', which requires an acknowledgement within 48~ms. Clients subsequently count timing events to maintain time.   Alternatively, PTP is being explored jointly by the Integrated Development Team (IDT) and the Ethernet to CAN (E2C) project \cite{e2cDescription}.  If PTP is demonstrated to be able to quickly ($<$1~sec) synchronize the time to within $\pm24$~ms (i.e. one half of a TE period), then client devices will be able to use PTP to maintain their local time. When ready to initialize an array of antennas, the control system then merely sends a command to ``align your clock to timing signals'', at which point the client devices round their local time to the nearest allowed TE value, and subsequently count timing events to maintain time.  In either case, once an initial time stamp is set referring to the TE when the control system starts (CSE: clock set event) all arrays will have the same time: the array time.  Every digitizer module must keep counting the 125~MHz reference and keep generating time stamps as long as the (sub)-array is active.  The minimum required accuracy of the time stamping in terms of absolute time is set by the VLBI observing mode at 1~microsecond, with a goal of 100~ns for the pulsar timing use case~\cite{sys_reqs}.


%%%%%%%%%%%%%%%%%%%%%%%%%%%%%%%%%%%%%%%%%%%%%%%%%
\subsubsection{Data Transmission System (DTS)}
%%%%%%%%%%%%%%%%%%%%%%%%%%%%%%%%%%%%%%%%%%%%%%%%%
\label{sec:DTS}
%%%%%%%%%%%%%%%%%%%%%%%%%%%%%%%%%%%%%%%%%%%%%%%%%
\paragraph{DTS Requirements and Constraints}
%%%%%%%%%%%%%%%%%%%%%%%%%%%%%%%%%%%%%%%%%%%%%%%%%
\label{DTSconstraints}
As noted in Section~\ref{digitizerSection}, the digitization process is expected to produce data sampled at 40~GSps per sideband per polarization with a bit-depth of 6. As a result, the required data rate will be 6b $\times$ 2pol $\times$ 2SB $\times$ 40~GSps = 960~Gbps. Added to this base rate are the required overheads, including the digitizer data time stamps, the antenna/sideband/polarization origin codes, and the Ethernet frame packing efficiency (total for all effects is $\approx$12\%,~\cite{ATAC_DSPOT_EICD}).\footnote{The data error correction (``Forward Error Correction'', or FEC) amounts to only $\sim$2\%, but
this additional protocol is supported over and above the individual link rate budget of 400~GbE.} Thus, a comfortable realistic placeholder target rate of 1.2~Tb/s per antenna will be adopted as a minimum requirement commensurate with the 6~bits/sample recommendation made in~\cite{scwg}. To avoid additional technical complexity to map the four digital data channels at each antenna, 2~polarizations / 2~sidebands, onto only three 400~GbE channels, it has been decided to allocate one 400~GbE to each sideband/polarization channel. This design solution results in a DTS with 4$\times$400~GbE optical channels, with a total capacity of 1.6~Tb/s, thus also meeting the recommended stretch goal of 8~bits/sample.\footnote{As a result, for the WSU use case, only $\approx$68\% of the continuously transmitted 400~GbE data will be properly formatted and meaningful while the rest will effectively be discarded at the receiving end.} This 400~GbE rate, as well as the other major key specifications for the DTS envisioned to meet the requirements of the WSU (see~\cite{CD_DTS}), are summarized in Table~\ref{tab:DTSreqs} alongside the corresponding properties of the current DTS (see~\cite{BEreqs}) to illustrate the large increase in performance.

\begin{center}
\begin{talltblr}
   [caption={Data Transmission System (DTS\index{Data Transmission System}) properties}, 
   label={tab:DTSreqs},
   remark{Notes}={Column 2 lists the current DTS properties, and columns 3 and 4 list the requirements to support 4$\times$ bandwidth and extended baselines.\\$^a$Subsystem requirements for DTS have not yet been established.  70 is currently used in this document as a placeholder value that can accommodate some future expansion in number of antennas.},
   ]
   {
    cells = {font=\small},
    colspec={|l|l|l|l|},
    row{1} = {headercolor},
    stretch=0.6pt,
   }
\hline
\bf Parameter&	\bf Legacy design	& \bf ALMA 2030 design & \bf Comment\\
\hline
Number of bits	& 3	& 6	& 6b\,+\,6b complex\\
Data rate	& 120 Gb/s	& 1.2 Tb/s	&  \\ %Several \\% overhead needed for FEC
Maximum Span &	15 km &	80 km	& 50 km to AOS + 30 km to OSF\\
Modulation format &	OOK/AM	& DP-16QAM/Coherent	& \\
WDM usage &	Yes: 12$\times$10 Gb/s & Yes: 4$\times$400 Gb/s	& \\
Bit Error Rate (BER) &	$10^{-6}$ &	$10^{-15}$ &	Forward Error Correction enabled\\
Protocol &	SONET/SDH &	Gigabit Ethernet	& \\
3R usage &	No	& No &	Optical amplification only\\
Number of antennas &	66	& 70$^a$ & \\
Fibers per antenna &  1  & 2  & Required for 400G COTS\\
\hline
\end{talltblr}

\end{center}

%%%%%%%%%%%%%%%%%%%%
\begin{comment}
\begin{table}
\small
\caption{Current Data Transmission System properties (column 2) and WSU requirements to support 4$\times$BW (columns 3 and 4).\label{DTSreqs}}
\setlength\tabcolsep{5pt}
\begin{tabular}{|l|l|l|l|}
\hline
\bf Parameter&	\bf Legacy design	& \bf ALMA 2030 design & \bf Comment\\
\hline
Number of bits	& 3	& 6	& 6b+6b complex\\
Data rate	& 120 Gb/s	& 1.2 Tb/s	& Several \\% overhead needed for FEC
Maximum Span &	15 km &	80 km	& 50 km to AOS + 30 km to OSF\\
Modulation format &	OOK/AM	& DP-16QAM/Coherent	& \\
WDM usage &	Yes: 12$\times$10 Gb/s & Yes: 4$\times$400 Gb/s	& \\
Bit Error Rate (BER) &	1e-6 &	1e-15 &	Forward Error Correction enabled\\
Protocol &	SONET/SDH &	Gigabit Ethernet	& \\
3R usage &	No	& No &	Optical amplification only\\
Number of antenna &	66	& 70 & \\	
\hline
\end{tabular}
\end{table}
\end{comment}
%%%%%%%%%%%%%%%%%%%%


In addition to these key specifications, operational concerns and constraints must be met to minimize the array downtime for science operations during the WSU upgrade. In the concept of parallel deployment (Section~\ref{sec:parallelDeployment}), the current 120~Gb/s system must continue to transmit during the WSU transition period. As a result, the antenna cabins, the Heating, Ventilation and Air Conditioning (HVAC) system, and the digital racks will have to be modified, so that the new Digitizers and DTS can be installed, powered up, tested, and operated with minimal disruption to the legacy system. This arrangement will require additional mounting points and hardware on the side panels of the FEs and digital racks for example (Fig.~\ref{fig541}), and additional cabling for the DTS signal, power distribution, and M\&C. The feasibility and impact of these changes in the antenna cabin, including RFI concerns, has been assessed by a joint working group %\sout{(see section 6.2)} 
\cite{IFSwitchLocation} and their recommendations on the possible options are informing the ongoing WSU design process in this area.

\begin{figure}
\includegraphics[width=0.40\textwidth]{Figures/AECcabin.jpg}
\includegraphics[width=0.59\textwidth]{Figures/Floorplan.png}
\caption{Left) The receiver cabin of an ALMA 12-m DA antenna viewed from the cabin entrance (credit: Sergio Otarola); Right) floor plan of receiver cabin racks in the Melco antenna~\cite{MelcoAntennaCabin}.
\label{fig541}} 
\end{figure}

%%%%%%%%%%%%%%%%%%%%%%%%%%%%%%%%%%%%%%%%%%%%%%%%%
\paragraph{Availability of Data Center Interconnect (DCI) architecture}
%%%%%%%%%%%%%%%%%%%%%%%%%%%%%%%%%%%%%%%%%%%%%%%%%
To mitigate the long-range effects of pulse spreading due to Chromatic Dispersion and Polarization Mode Dispersion, the current Direct Detection Amplitude Modulation (OOK) method widely used in lower data rate systems (40~Gb/s or less) is being replaced by Coherent modulation (ex: QPSK: Quadrature Phase Shift Keying). High performance, long range 100~Gb/s, 200~Gb/s, and 400~Gb/s are all Coherent type formats. The recommended ``Physical Layer'' data transmission hardware family that can meet the WSU target specs is the current industry-developed family of 400~GbE CFP2/4 or QSFP-DD Digital Coherent Optics (DCO) Transceivers. This choice would allow the entire data stream from each antenna to be carried by four 400~Gb/s transmitters operating in the optical C-Band (1550~nm). This wavelength corresponds to the minimum attenuation window of the existing ALMA optical fiber network and it allows the (re)use of telecom grade WDM components and EDFAs.

%%%%%%%%%%%%%%%%%%%%%%%%%%%%%%%%%%%%%%%%%%%%%%%%%
\paragraph{DTS Interface to Digitizer FPGA}
%%%%%%%%%%%%%%%%%%%%%%%%%%%%%%%%%%%%%%%%%%%%%%%%%

As described in the digitizer prototyping project that will build upon the work in the Digitizer upgrade study~\cite{Quert_final}, the digitizer outputs will be captured and processed by an FPGA stage, the DSP/OT, likely co-located near the FE or BE and DTS electronics. The signals will be transmitted via four 400~GbE optical transceivers (see Fig.~\ref{fig543}).  The interface of these optical transceivers to the DSP/OT will be an industry standard being either QSFP-DD or QSFP-112. The resulting four optical signals must then be multiplexed onto one of the selected spare single-mode fibers using a commercial DWDM device. Further details and options are described in the DTS project proposal~\cite{CD_DTS}.

%%%%%%%%%%%%%%%%%
\begin{figure}
\includegraphics[width=1.0\textwidth]{Figures/MihoFig543.png}
\caption{Functional schematic of the digital data flow from the optical transceivers in each antenna to the AOS building and to the OSF building, which will house the new correlator (Section~\ref{sec:ATAC}) and the TPGS (Section~\ref{sec:TPGS}).
\label{fig543}} 
\end{figure}
%%%%%%%%%%%%%%%%%


%%%%%%%%%%%%%%%%%%%%%%%%%%%%%%%%%%%%%%%%%%%%%%%%%
\paragraph{DTS Interface to AOS Patch Panel Room}
%%%%%%%%%%%%%%%%%%%%%%%%%%%%%%%%%%%%%%%%%%%%%%%%%

Since the proposed QSFP-DD DCO modules have relatively low optical output power, and the new correlator will be located at the OSF, an additional 30~km away from the AOS, simple COTS inline EDFAs must be installed in the Antennas and the Patch Panel rooms of the AOS and OSF to boost the signals (to allow for margin and ageing). See Section~\ref{WSU_EDFA_rack} for further details.

%%%%%%%%%%%%%%%%%%%%%%%%%%%%%%%%%%%%%%%%%%%%%%%%%
\paragraph{DTS Interface to New Correlator (and ACA Spectrometer)}
%%%%%%%%%%%%%%%%%%%%%%%%%%%%%%%%%%%%%%%%%%%%%%%%%
The ATAC (Section~\ref{sec:ATAC}) equipment can ingest up to 1200~Gb/s from each antenna. The latest design of the ATAC correlator is based on Intel\,\textsuperscript{\textregistered} Agilex FPGAs, which are capable of handling 3$\times$400~GbE, 6$\times$200~GbE, or 12$\times$100~GbE data rates. Furthermore, a 400~GbE Ethernet switch will be installed between DTS and the ATAC correlator FPGA ingest. The latest Ethernet switch can also directly ingest 400~GbE inputs for long reach optics and provide power supply up to 24~W for the optics. Thus, the 400~GbE inputs on the switch shall be the physical interface between DTS and the ATAC correlator subsystem. Additionally, the Ethernet switch will be used to forward the data packets from the four TP antennas to the upgraded ACA Spectrometer (Section~\ref{sec:TPGS}).


%%%%%%%%%%%%%%%%%%%%%%%%%%%%%%%%%%%%%%%%%%%%%%%%%
%DTS Power Requirements}
%%%%%%%%%%%%%%%%%%%%%%%%%%%%%%%%%%%%%%%%%%%%%%%%%
%For the proposed DTS design, a rough estimate of the power requirements for the three major assemblies (70 antennas, AOS Patch Panel Room, Correlator Room) can be estimated as follows:
%\begin{enumerate}[itemsep=-0.5em]
%\item %Antenna: 4 $\times$ QSFP-DD-DCO modules + EDFA = 120~W per antenna.
%\item %AOS Room: 20 EDFAs + PDU + M\&C = 600~W.
%\item %Correlator Room: 280 $\times$ QSFP-DD-DCO modules + 70 $\times$ EDFAs+ PDU + M\&C = 10~kW total for 2 racks to support 70 antennas total.
%\end{enumerate}
%As such, as placeholder estimate for the power consumption of the new hardware is approximately 19~kW (70~antennas). The uncertainty is 20\%, thus we should allow 150~W per antenna in the power budget (see Table~\ref{tab:powerBudget}) until the design is more mature.  %The equipment in the antenna does not need to be on a UPS unless there is a reason that the existing DTS is already on a UPS.

%\begin{center}
\begin{talltblr}[caption={DTS\index{Data Transmission System} Electrical power consumption estimation: TO BE UPDATED BY CARLA},label={tab:powerBudget}]
{
colspec={QQQQ},
cells={font=\small},
row{1}={headercolor, font=\small\bfseries},
row{2,5}={highlight2},
hlines,
vlines,
}
Component & kW & kVA & UPS needed?\\
\SetCell[c=4]{c} AOS\\
DTS\index{Data Transmission System} per-antenna component & 0.2 & 0.2 & No\\
DTS\index{Data Transmission System}-AOS component & 0.6 & 0.7 & No\\
\SetCell[c=4]{c} OSF\\
DTS\index{Data Transmission System}-OSF component & 10 & 11 & No
\end{talltblr}
\end{center}

%%%%%%%%%%%%%%%%%%%%%%%%%%%%%%%%%%%%%%%%%%%%%%%%%
\subsubsection{Advanced Technology ALMA Correlator (and Beamformer)}
%%%%%%%%%%%%%%%%%%%%%%%%%%%%%%%%%%%%%%%%%%%%%%%%%
\label{sec:ATAC}

The ALMA science goals, updated for the WSU~\cite{road_map} and their flow-down to specifications for the 2nd Generation Correlator~\cite{2ndgen}, and subsequently the WSU System Requirements~\cite{sys_reqs}, require:
%%%%%%%%%%%%%%
\begin{itemize}[itemsep=-0.5em]
    \item A huge increase in the total number of channels available.\footnote{It must be possible to achieve a spectral resolution of $\sim$0.1 km/s at any ALMA frequency from 35 to 950~GHz at full correlated bandwidth, a feat that cannot now be achieved for any ALMA band. For reference at 35~GHz, 13.5~kHz $\sim$0.1 km/s and to cover 16~GHz of bandwidth requires 1.2~million channels per polarization product.}
    \item A significant increase in correlation bit depth relative to the existing ALMA correlators, as well as
    \item The ability to do either correlation or beamforming across the full bandwidth, which should be at least a factor of two larger than the current correlators, with a goal of four times larger.
\end{itemize}
%%%%%%%%%%%%%%
Additionally, it should meet these goals without significantly increasing the total power (hardware + coolling) consumption required, compared to the current correlators. Fulfilling these requirements is a major engineering effort that must fit within the funding and resource constraints of the ALMA development program. 
To address these challenges, ATAC has chosen the Intel Agilex-M\,\textsuperscript{\textregistered} family of FPGAs\footnote{https://www.intel.com/content/www/us/en/products/details/fpga/agilex/7/m-series.html} as the workhorse of the 2nd Generation ALMA correlator. These FPGAs have many advantages over the previous generations of FPGAs in terms of lower power consumption (for the same processing job), higher performance fabric to help meet timing, greater FPGA resources per device (fewer devices needed), faster transceivers and better hard IP for communications, and higher memory bandwidth (both DDR5 and High Bandwidth Memory available). Additionally, FPGAs are uniquely suited to the ALMA 66-antenna, wide band with use case, because the same hardware resources can be optimized for the diverse needs required for:
%%%%%%%%%%%%%%%%%%%%%%%%%%%%%%%%
\begin{itemize}[itemsep=-0.5em]
\item Imaging Correlation and fine channelization (i.e. standard interferometry)
\item VLBI Beamforming with associated coarse channelization visibilities
\item Single antenna VLBI capability (without affecting the ability to use the rest of the array elements)
\item Hardware in the loop capability for 4-antennas at full science bandwidth (using a distinct software version and without affecting the use of other array elements).
\end{itemize}
%%%%%%%%%%%%%%%%%%%%%%%%%%%%%%%%

The Advanced Technology ALMA Correlator (ATAC) will meet these goals, and provide the following capabilities:

%%%%%%%%%%%%%%%%%%%%%%%%%%%%%%%%%%
\begin{enumerate}[itemsep=-0.5em]
%%%%%
\item	Ingest digitized signals from up to 66 antenna inputs, each with a bandwidth of 16~GHz per sideband per polarization (i.e. 4$\times$BW, digitized at 40.0~GS/s), provided that the total data rate (including overhead) is $<$1.2~Tbps (see Fig.~\ref{fig:ATACFig1} for a functional overview).
%%%%
\item Capacity to correlate or beamform up to 32~GHz of bandwidth per polarization (4$\times$BW) when fully deployed.  The initial deployment of ATAC (i.e. the currently funded project) will support 16~GHz of correlated or beamformed bandwidth per polarization (i.e. 2$\times$BW).  The 16~GHz of initially deployed bandwidth can be selected from either sideband in 200~MHz increments (all 16~GHz of bandwidth can come from a single sideband or any combination of bandwidth, in 200~MHz increments, selected from both sidebands). These 200~MHz Frequency Slices (FSs) are the building blocks of flexible, user-defined spws that are formed downstream in the Correlator Data Processor (CDP).
%%%
\item	Processing includes delay compensation, LO offset removal, and Earth-rotation phase correction. All signal processing occurs at a bit-depth $\geq$8-bits, apart from the final correlation stage where the bit-depth is 6-bits, for a digital sensitivity improvement in the correlator of 13.4\% compared to BLC FDM modes. 
%%%%%
\item	Support for at least six correlator subarrays.
%%%%%
\item In IC mode, produce a native fine channel bandwidth of 13.5~kHz at full correlated bandwidth (even finer channels are available by factors of 2, 4, 8, and 16, with proportional decreases in correlated bandwidth), producing up to $\sim$1.2 million channels with a time cadence as small as 1~second, for the 2$\times$ IF bandwidth upgrade. At 80$\times$ coarser channel bandwidth (1.08~MHz), a time cadence as fine as 16~ms is available.\footnote{Coarser channel bandwidth can always be obtained by selecting integer values of channel averaging.} All four polarization products are always produced (though the cross-hands can be dropped if desired) Flexible time averaging and channel averaging will be possible on a per spectral window basis.
%%%%%
\item	In VLBI Beamforming (VB) mode, for each selected frequency slice, form one VLBI voltage beam at a chosen delay center on the sky.  Processing uses true-delay beamforming to generate each beam.\footnote{After application of WVR-corrected phasing solutions.}  In addition to VLBI beams, visibilities with relatively coarse channel bandwidth (218.75~kHz) are simultaneously generated, processed, and formatted in a similar manner to IC visibilities, and used in near-real-time for computing phasing solutions.
%%%%%
\item	A stand-alone single-antenna VLBI mode (A1-VLBI) where one of two pre-defined antennas (configured at the AOS patch panel) can be selected to convert the wideband digitized data into VLBI data streams, which are transmitted to external VLBI recording equipment.  Processing includes coarse channelization, resampling to VLBI-``friendly'' sample rates, and packetization / transmission to external VLBI equipment. 
%%%%
\item	A stand-alone slice of ATAC hardware for use in the WSU Hardware in the Loop Simulation Environment (WSU HILSE, Section~\ref{sec:WSU_HILSE}) to perform tests using either sky signals from four ALMA antennas or four simulated digital signals generated by WSU HILSE equipment and up to 2~GHz of processed bandwidth per polarization.
\end{enumerate}
%%%%%%%%%%%%%%%%%%%%


%%%%%%%%%%%%%%%%%%%%
\begin{figure}[htb]
\includegraphics[width=\textwidth]{Figures/ATACFig1.png}
\caption{Overview of ATAC subsystem showing the primary functionality. The red dashed box shows the expansion path to 4$\times$BW Correlation \& VLBI Beamforming. \label{fig:ATACFig1}} 
\end{figure}
%%%%%%%%%%%%%%%%%%%%


ATAC achieves these capabilities with a two-stage FFX architecture, followed by a final Correlator Data Processor (CDP) stage (also see Fig.~\ref{fig:ATACFig1}): 

\textbf{Stage 1: The ALMA Very Coarse Channelizers (AVCCs)} are antenna-based and perform the first ``F'' channelization that is common to the IC, VB and single-antenna VLBI (A1-VLBI) observing modes, producing oversampled ``Frequency Slices'' (FSs) via an oversampling poly-phase filterbank (OSPFB).  These FSs are 222.222~MS/s wide (for each polarization), with the central (slightly more\footnote{Sufficient to accommodate LO offsetting, earth-rotation phase, and offsets to facilitate stitching FSs.} than) 200~MHz containing un-aliased science-quality data.  Each AVCC generates 81 dual-polarization FSs\footnote{An extra FS is required for each sideband to handle the case where the edges of the science bandwidth within the wideband digitized signal does not align with the FS boundaries (100~MHz, 300~MHz, etc.).} (162~total) to cover the required 16~GHz.  The 16~GHz extent covered by the 81~FSs per sideband can be selected from anywhere in the 20~GHz of digitized bandwidth to support receivers for different bands that produce science bandwidth in a different region of the digitized bandwidth.  Each FS can be transmitted to one or more ALMA Frequency Slice Processors (AFSP) to perform further processing.  

\textbf{Stage 2: The ALMA Frequency Slice Processors (AFSPs)} are Frequency Slice based and can be configured to run in one of three modes: Imaging Correlation, VLBI Beamforming, and A1-VLBI.  Each AFSP can only operate in one mode at a time but can be reconfigured at will to match the needs of a particular observation.  Each AFSP also performs the antenna-based pre-processing (delay tracking, LO offset removal and fine channelization via a critically sampled poly-phase filterbank) on one 200~MHz FS for all 66 ALMA antennas.  Each AFSP can perform its selected mode (cross correlation or beamforming) for all correlator subarrays, even if each correlator subarray has different channelization and integration time parameters.  See Fig.~\ref{fig:ATACFig2} for an overview of how the different stages of channelization relate to each other.

%%%%%%%%%%%%%%%%%%%%
\begin{figure}[htb]
% going wider than 0.92 kicks it to end of section
\includegraphics[width=0.92\textwidth]{Figures/ATACFig2.png}
\caption{Visualization of the IF range options for upgraded receiver bands and subsequent channelization stages of ATAC for Imaging Correlation mode. The AVCC stage only illustrates the FSs from one sideband, though both sidebands are channelized to 200~MHz FSs. For subsequent processing, for the initial 2$\times$BW correlation or beamforming deployment, up to 80~FSs can be selected for AFSP processing selected flexibly from either sideband. Frequency-continuous FSs can be stitched into wider user specified spws.  The light blue IF range in the top ``antenna'' panel is indicative of the IF ranges planned for Bands~6v2 and 8v2, while the teal colored IF range is indicative of Band~2. \label{fig:ATACFig2}} 
\end{figure}
%%%%%%%%%%%%%%%%%%%%


For example, if all 81 AFSPs in the initial 2$\times$BW deployment are configured for Imaging Correlation, one correlator subarray of $N$ antennas could be using all 81 AFSPs to produce a single 16~GHz spw with wide channels ($\geq$1.08\,MHz) and short integration times (16~ms).  Simultaneously, another correlator subarray of $66-N$ antennas could be using 41 AFSPs for an 8~GHz spw and another 40 AFSPs for zoom windows on spectral lines, all with longer integration times ($\sim$3 second).  In another example, if one correlator subarray is using 20 AFSPs to perform VLBI Beamforming on a 4~GHz spw, there would only be 61 AFSPs available for another correlator subarray to use for Imaging Correlation. ATAC is designed so that 4$\times$BW correlation or VLBI/Beamforming can be achieved with a straightforward future hardware expansion of only the Stage~2 (AFSP) processing hardware.

\textbf{Stage 3: The Correlator Data Processor (CDP)} receives visibility sets from AFSPs operating in Imaging Correlation or VLBI Beamforming (including low spectral resolution visibilities to support VLBI beamforming) mode.  The CDP will seamlessly combine multiple, adjacent 200~MHz FSs together to cover the contiguous bandwidth requested by the user for each spw. For each spw, the CDP performs any additional channel and time averaging of visibility sets generated by the AFSP to meet the specified spectral resolution and integration time.  The CDP also performs the Van Vleck corrections on visibilities, normalizes cross-correlation visibilities based on power spectra of the corresponding auto-correlations (to flatten the bandpass), generates BDFs, and transmits them to TelCal\index{TelCal} and Archive. An innovative aspect of the ATAC design is that the bulk of CDP processing will occur on CPUs that are co-located with the AVCC and AFSP hardware (aside from the CDP Manager and CDP Formatter Servers).


%%%%%%%%%%%%%%%%%%%%%%%%%%%%%%%%%%%%%%%%%%%%%%%%%
\subsubsection{Total Power GPU Spectrometer (TPGS)}
%%%%%%%%%%%%%%%%%%%%%%%%%%%%%%%%%%%%%%%%%%%%%%%%%
\label{sec:TPGS}
The Atacama Compact Array Spectrometer (ACA Spectrometer), a spectrometer for the TP Array, based on GPU technology, was installed at the AOS in early 2022, and has been used for PI science observations from the start of Cycle~10. In the coming era of the WSU, the ACA Spectrometer will be succeeded by an advanced subsystem, the TPGS. This transition will make use of the foundational SW of the ACA Spectrometer to facilitate the development of the TPGS.

Recently, GPU technology has grown rapidly due to the increasing demands of artificial intelligence applications. This progress has significantly increased the computing power of GPUs and enhanced memory access bandwidth. A landmark in this evolution is the advent of PCIe 5 and the NVIDIA Grace Hopper superchip, which significantly boost data transfer rates from the antennas to the GPU memory and provide superior processing capabilities. The integration of such state-of-the-art GPU technology into the TPGS will be necessary to meet the demanding WSU requirements, including increased data rates from the TP Array, much finer spectral resolutions, and many spws.

Designed to generate full polarization products from digitized samples for each of the four TP antennas, the TPGS will also support the pointing and focus calibrations for the TP Array by providing cross-correlation products with optimized bandwidths and spectral resolutions. The TPGS will employ a FFX architecture, processing sampled data on a spw basis. It will accommodate digitized samples across a 20 GHz bandwidth per polarization per sideband per antenna, utilizing a polyphase filter to isolate spws specified by the PI within the 20 GHz range. Subsequently, it will perform the FFT on the filtered samples to compute the auto-correlation. For full polarization products, the TPGS will select identical spws from digitized sample streams of X and Y polarizations, computing both auto-correlation and cross-polarization correlation. Its design will allow flexibility in selecting the bandwidth of a spw and the FFT size, thereby fulfilling the requirements for different bandwidths of spws and their spectral resolutions.

The TPGS subsystem will be installed within a rack in the newly established OSF correlator room, sharing 400~GbE switches with the ATAC. It will be tailored to separately process USB and LSB data streams. For instance, all USB data streams from the four TP antennas will be routed through the 400 GE switches to a server equipped with multiple GPUs. Additionally, the TPGS HILSE will be co-located on the same rack, designed to process one sideband data stream from four antennas.


%%%%%%%%%%%%%%%%%%%%%%%%%%%%%%%%%%%%%%%%%%%%%%%%%
\subsubsection{Site Infrastructure}
%%%%%%%%%%%%%%%%%%%%%%%%%%%%%%%%%%%%%%%%%%%%%%%%%
As described in \cite{2ndgen} and \cite{scwg}, there are many reasons to install the new ALMA correlator at the OSF rather than the AOS. The implications of this relocation are discussed in the following subsections, followed by the electrical power budget needed for the WSU. 


%%%%%%%%%%%%%%%%%%%%%%%%%%%%%%%%%%%%%%%%%%%%%%%%%
\paragraph{New AOS to OSF Fiber Optic Cable}
%%%%%%%%%%%%%%%%%%%%%%%%%%%%%%%%%%%%%%%%%%%%%%%%%
\label{sec:newFiber}
As noted in~\cite{scwg}, the new fiber cable must provide the dedicated fiber paths for the digitized signals from each antenna to the new correlator. Considering that there are currently 66 antennas, allowing for a moderate expansion in number of antennas (up to 80, which is the current limit of the available CLO hardware: 66 + 15\% spares + 4 for HILSE) and allowing for a spare fiber for each antenna we need at least $80 \times 2=160$ fibers. As described in~\cite{optcab}, the fiber count available in a typical cable is 48, 64, 96, 144, or 192 cores.  In the initial study led by G. Filippi in 2021, options have been costed for 192 fibers (in two cables of 96 each) or 3$\times$96 fibers (in three cables of 48 pairs each). It is expected that the optical cable specification for the new cable will be identical to the current one~\cite{befo}.
 
The functionality and usage of the existing 48-core AOS-OSF communications cable will not be changed during the WSU commissioning period, as it includes the flow of data from the BL Correlator and ACA Spectrometer to the archive.\footnote{The ACA functionality for the 7-m Array will be replaced by the BLC~\cite{ACAEoL}. The ACA functionality for the TP Array was replaced by the ACASPEC in Cycle 10.} Of these 48 fibers, 27 are currently available as defined in \cite{FO_Memo}. However, eighteen fibers will need to be enabled or extended for the initial part of the WSU deployment, as they are currently only available between the power plants at the AOS and OSF, but not at the technical buildings. The distances between the power plants and the technical buildings are between 150~m and 400~m.



%%%%%%%%%%%%%%%%%%%%%%%%%%%%%%%%%%%%%%%%%%%%%%%%%
\paragraph{Modifications to the AOS Fiber Patch Panel}\label{AOS_patchpanel}
%%%%%%%%%%%%%%%%%%%%%%%%%%%%%%%%%%%%%%%%%%%%%%%%%
The current four-rack fiber patch panel at the AOS (Fig.~\ref{fig:AOSpatchPanel}) accepts the 8 fibers from each antenna pad (DTS\,+\,LO\,+\,2~M\&C\,+\,4 spares) on the rear side. The specification for the fibers in the AOS building is given in~\cite{befo}.

On the front side, the 2$\times$192 M\&C fibers are labelled by pad name and are always connected to all 192 pads. In contrast, the 66 DTS and 66 LO fibers are labelled by antenna name and are connected manually to the sockets corresponding to the currently occupied pads. The 66 emerging DTS fibers are gathered into a fiber cable that is routed to the two correlator rooms (one for the BLC, and one for the ACAC/ACAS). When an antenna is moved to a new pad, its DTS and LO fiber connections are moved manually from the old pad’s connector module to the new pad’s connector module~\cite{PatchPanel}. The parallel deployment philosophy of the WSU means that this patch panel must remain in operation and not be impacted by the installation and testing of the WSU. However, a whole separate patch panel is not necessary, because the spare fibers from each antenna are already terminated in the existing patch panel. Therefore, the minimum required change to the patch panel is the addition of a new short ($\sim$10~m) fiber cable that gathers the spare fiber connections from the currently-selected antennas on the existing patch panel and routes them to a new single rack (Section~\ref{WSU_EDFA_rack}).  

%%%%%%%%%%%%%%%%%%%%
\begin{figure}[htb]
\includegraphics[width=\textwidth]{Figures/AOSpatchPanel.jpg}
\caption{Photograph of the AOS patch panel, obtained in March 2022. \label{fig:AOSpatchPanel}} 
\end{figure}
%%%%%%%%%%%%%%%%%%%%


%%%%%%%%%%%%%%%%%%%%%%%%%%%%%%%%%%%%%%%%%%%%%%%%%
\paragraph{Addition of a WSU/EDFA Rack}
%%%%%%%%%%%%%%%%%%%%%%%%%%%%%%%%%%%%%%%%%%%%%%%%%
\label{WSU_EDFA_rack}
This new ``WSU/EDFA rack'' will contain about 20 inline EDFAs that can be used to boost the optical power of the signal from any antennas that are occupying a distant pad with power below the required threshold to reach the OSF. The current photo of the patch panel room (Fig.~\ref{fig:AOSpatchPanelRoom}) shows that there is plenty of space available for this new rack. The spot currently planned is the empty rack slot with a seismic mount on the left of the  Fibre Optic Amplifier Device (FOAD) rack.

%%%%%%%%%%%%%%%%%%%%
\begin{figure}[htb]
\includegraphics[width=\textwidth]{Figures/AOSpatchPanelRoom.jpg}
\caption{Patch panel room showing the available empty rack mounting slot (right of center). \label{fig:AOSpatchPanelRoom}} 
\end{figure}
%%%%%%%%%%%%%%%%%%%%

The output from this new WSU rack will then go through another new short fiber ($\sim$30~m) cable to the room where the new AOS-OSF fiber (Section~\ref{sec:newFiber}) is to be terminated.  These two new short fiber cables should contain at least 2$\times$92 fibers, as there will be 2 fibers per antenna (Section~\ref{sec:DTS}) and 80 is the number of antennas that the currently-installed CLO hardware supports, so allotting 15\% spares requires 12 more. During the WSU commissioning phase, when an antenna is moved to a new pad, then four of its fiber connections (LO\,+\,DTS\,+\,WSU DTS) will need to be moved instead of the (only) two fiber connections that need to be moved currently (LO\,+\,DTS).


%%%%%%%%%%%%%%%%%%%%%%%%%%%%%%%%%%%%%%%%%%%%%%%%%
\paragraph{OSF Correlator Room (OCRO)}
%%%%%%%%%%%%%%%%%%%%%%%%%%%%%%%%%%%%%%%%%%%%%%%%%
\label{sec:OCRO}
The current control room at the OSF will be refurbished to house and support the specific needs of the new ATAC correlator (Section~\ref{sec:ATAC}), the TPGS (Section~\ref{sec:TPGS}), and the new DTS fiber receiving electronics (Section~\ref{sec:DTS}). The refurbishment will require establishing a room with appropriate entry protection, adding seismic protection for the racks, providing sufficient airflow and HVAC capacity for cooling the equipment, providing a fire protection system, and connecting to the electrical power system. The room should be designed considering the running cost of cooling, as well as the capital cost, using efficient cooling methods when outdoor conditions allow. The increase in power draw for this room means that the existing transformers need to be assessed for their capacity, and appropriate steps taken if it is insufficient in any respect (see Section~\ref{sec:powerBudget}). The current control room has ample space for the upgraded correlator and associated components, as shown in Fig.~\ref{fig:OCROlayout}  (28 racks in 2 rows: 2 racks for DTS, 2 for TPGS, 22 for a complete 4$\times$BW ATAC, 4 spare racks allocated to future computing needs, and 2 liquid cooling heat exchangers). 

A new location for the OSF control room has been identified, as it is still used regularly for the observing night shift. Construction work for the Control Room Relocation (CRR) project is expected to start in Q2 2024. The size and functionality of the new control room will be similar to the SCO control room.


%%%%%%%%%%%%%%%%
\begin{figure}[t]
\includegraphics[width=\textwidth]{Figures/Drawing_for_WSU_CoSDD_Feb2024.png}
\caption{Left panel) Conceptual diagram of the existing OSF control room, with its entrance from the main hallway of the OSF technical building. Right panel) Draft layout of the room after refurbishment to house the new ATAC correlator and TPGS spectrometer, including rack space (600~mm wide units with 19-inch internal mounting) for the future 4$\times$ bandwidth configuration.  The DTS equipment will be located centrally to minimize the lengths of subsequent connections to the ATAC and TPGS racks. Ample space for the new AOS-to-OSF fiber breakout panel is available in the upper right area. \label{fig:OCROlayout}} 
\end{figure}
%%%%%%%%%%%%%%%%


%%%%%%%%%%%%%%%%%%%%%%%%%%%%%%%%%%%%%%%%%%%%%%%%%
\paragraph{Electrical Power Budget and Power Generation}
%%%%%%%%%%%%%%%%%%%%%%%%%%%%%%%%%%%%%%%%%%%%%%%%%
\label{sec:powerBudget}
On the time frame of the WSU deployment, the extension of the national power grid to reach ALMA has low prospects. Alternative power sources may be sought to augment the existing turbine power (e.g. renewable like a solar panel farm %with battery storage backed the existing generator set
) to provide additional power support to run some portions of the OSF, such as the Residencia (500~kW), freeing more turbine power to support the parallel deployment of the WSU hardware. The costs, for both construction and operation phases, of these alternative power sources need to be carefully balanced against the costs of designing more energy efficient equipment. 

Currently, ALMA has 3 turbines, with a maximum operational value of 3.3~MW per unit. From the 3 turbines, one of them operates and the other 2 are in standby mode. For stability of the power generation system, the average power demand must be below 3~MW for the turbine in operation.
The monthly current average power demand (without WSU) is between 2.4~MW and 2.5~MW.
For the WSU parallel deployment phase, the preliminary estimation of the additional power budget for the new WSU elements is 761~kW for the OSF equipment and 1.2~kW per each antenna retrofitted for WSU, during the transition to WSU. 
Mitigation studies are needed for the supply of the additional electrical power required during the parallel deployment stage, if the intention is to avoid the use of a second turbine (as operating the second turbine for a few extra KW is very inefficient and not preferred).

A first version of the required additional Power Budget at a conceptual level (considering the maximum consumption possible during the parallel deployment phase) can be found in Table~\ref{tab:powerBudget}.


%\begin{spacing}{0.1}
\begin{center}
\begin{talltblr}[caption={WSU Additional electrical power consumption estimation},label={tab:powerBudget}]
{
colspec={Q[7cm]Q[c]Q[c]Q[c]Q[6cm]},
cells={font=\scriptsize},
row{1,2}={headercolor, font=\scriptsize\bfseries},
row{3,12}={highlight2},
row{10,11,20,22}={light_green, font=\scriptsize\bfseries},
hlines,
vlines,
width=\textwidth,
stretch=0.75pt,
}
\SetCell[r=2]{c} Component & \SetCell[c=2]{c} {Average\\demand} & &  \SetCell[r=2]{c}{UPS\\needed?} & \SetCell[r=2]{c} Assumptions\\
& kW & kVA & \\
\SetCell[c=5]{c} {\bf AOS}\\
Front-End additional instruments per antenna & 0.1 & 0.1 & No & Only the new Band 2 Rx will yield additional power consumption\\
WSU Back-End instruments per antenna & 0.28 & 0.28 & No & \\
WSU IF switch per antenna & 0.02 & 0.02 & No & Current IF switch will be replaced by upgraded version\\
Extra cooling for additional WSU Back-End per antenna & 0.0 & 0.0 & No & No need for additional cooling in antenna Rx cabin\\
DTS\index{Data Transmission System} per-antenna component & 0.2 & 0.2 & No \\
DTS\index{Data Transmission System}-AOS component & 0.6 & 0.7 & No  & \\
Total additional power per antenna & 1.2 & 1.3 & & \\
Total additional power at AOS for 66 antennas & 79.2 & 85.8 & & Worst case scenario 66 antennas, considers each antenna prepared for parallel deployment \\
\SetCell[c=5]{c} {\bf OSF}\\
ATAC\index{Advanced Technology ALMA Correlator} with 4x BW ingest and initial channelization and 2x BW Correlation/Beamforming.  & 228 & 253 & \makecell{Yes\\ (short \\ duration)} & Includes all equipment in ATAC racks for signal processing, Monitor \& Control, Correlator Data Processor compute, PDUs, in rack UPS for Monitor and Control, also includes in-rack liquid cooling power and the ATAC.HLSE. This value includes a 30\% margin.\\
HVAC for OSF Correlator Room  & 433 & 529 & No & To be refined after the HVAC study\\
%Additional HVAC for new correlator: option 3 & 210 & 263 & No & To be refined after the HVAC study\\
%HVAC Optional: liquid cooling for option 1 & 180 & 225 & No & To be refined after the HVAC study\\
%HVAC Optional: liquid cooling for option 3 & 210 & 263 & No & To be refined after the HVAC study\\
DTS\index{Data Transmission System}-OSF component & 10 & 11 & No & \\
TPGS\index{Total Power GPU Spectrometer} (Total Power GPU Spectrometer) & 40 & 48 & Yes & Power consumption estimated at the end of the upgrade: 40 kW\\
Software Servers Infrastructure & 22 & 25 & Yes & 5$\times$ Blade Chassis\\
Additional Archive capacity & 19.8 & 23.8 & Yes & 2 racks, worst case scenario \\
%Bulk-Data Server (Container de Archive) GAS 03 & 8 & 9 & Yes & 5$\times$ RS-840\\
IT equipment & 8 & 9 & Yes & 2 racks, worst case scenario \\
%Server APE (Software Group) & 8 & 9 & Yes & 3$\times$ Blade Chassis\\
Total additional power at OSF (ATAC\index{Advanced Technology ALMA Correlator} 2$\times$ bandwidth) & 760.8 & 898.8 & & Worst case for correlator \\
%Total additional power at OSF (option 3: 4$\times$ bandwidth respect to existing correlator - HVAC include) & 736  & 817 & & Worst case for correlator, with only air cooling\\
%\SetCell[c=5]{c} Total OSF plus AOS Correlator\\
%2$\times$ bandwidth & 599 & 666\\
%4$\times$ bandwidth & 806 & 896\\
Additional power consumption for the upgrade of the ATAC\index{Advanced Technology ALMA Correlator} to 4x BW Correlation/Beamforming. & 104 & 116 & TBC & The increase (after future upgrade) is for additional 2nd stage (AFSP) hardware only. This value includes a 30\% margin\\
Total additional power at OSF (ATAC\index{Advanced Technology ALMA Correlator} 4$\times$ bandwidth) & 864.8 & 1014.8 &   & \\
\end{talltblr}
\end{center}
%\end{spacing}

% Table 5.11.5 here


%%%%%%%%%%%%%%%%%%%%%%%%%%%%%%%%%%%%%%%%%%%%%%%%%
\subsubsection{Handling a Mix of Legacy and Upgraded Devices}
\label{handling_legacy_upgraded_devices}
%%%%%%%%%%%%%%%%%%%%%%%%%%%%%%%%%%%%%%%%%%%%%%%%%
As described in Section~\ref{sec:frontend}, not all receiver bands will be upgraded at once, so both the WSU hardware downstream of the receivers and the control SW will need to support a mix of legacy receiver bands and upgraded receiver bands. In addition, new hardware device drivers will be developed to control and/or monitor new hardware introduced in the WSU, such as the EDFAs to be added in the patch panel room to boost the optical signals to reach the OSF (Section~\ref{sec:DTS}).   The existing hardware device driver framework abstracts the Monitor and Control functionalities needed to interface with the hardware. This framework comprises a model to describe the Monitor and Control points (defined in the ICD) and a code generation to generate a common portion of the functionalities. SW developers could extend these base functionalities if it is necessary for a given hardware device. Currently, the hardware device framework supports drivers based on the CAN bus, RS-232, and Ethernet. For example, the Front End Monitor and Control (FEMC) module communicates via CAN bus, and the cryostat vacuum sensors communicate via RS-232.  Examples of Ethernet type of devices are the WeatherStation and the Central Variable Reference (CVR). The existing Ethernet device driver could be extended to support new protocols chosen for use by ALMA 2030 hardware. The challenge will be to choose protocol(s) that will persist for decades.
%such as MQTT, OPC/UA, or any other in-house protocol built on top of TCP/IP. 

The binaries generated by the device driver framework are run on the ABM, the main real-time computer located in each antenna to expose M\&C functionalities to the high-level observing SW. The hardware of the ABM is already obsolete. It will be replaced in the near future by the Ethernet to CAN (E2C)~\cite{e2cDescription} gateways, an embedded device solution based on an ARM processor along with integrated FPGA fabric, which will support the M\&C of legacy devices that won't be upgraded by the WSU. The E2C project has just passed the PDR review. It is expected the new hardware devices considered in the WSU will not use the CAN bus but more likely an Ethernet-based protocol. Therefore, E2C is not required by new WSU devices.  However, the next generation Front End Monitor Controller (ngFEMC), which is an obsolescence project, will have a CAN interface for two reasons: 1) so that the software switchover to using Ethernet communications in the antennas can be staged, and 2) so that FE test centers can connect to ngFEMCs without updating their software.

The changes introduced to support new devices in parallel with the existing ones in operation will impact the Telescope Monitor \& Control Database (TMCDB), APDM\index{ALMA Project Data Model}, ASDM\index{ALMA Science Data Model}, and others. These are important aspects of ALMA through which the various SW tools interact.  The TMCDB contains configuration schemas for the hardware devices, for which the new SW device drivers will require new entries. For example, the first LO’s warm multiplier harmonic value will be changed from 6 to 3 in the upgraded Band 6 WCA~\cite{b6up} (and possibly in other receiver bands), which will affect the FE code (but not the FLOOG).  

During the transition period, when upgraded receiver cartridges are gradually deployed, some antennas will have upgraded receivers for a given Band while others will still have legacy receivers for that Band, and the difference between antennas will need to be accommodated.  One option will be to make use of the existing capability that each device in the antenna can be launched with a different implementation of the corresponding SW driver (for example, under real mode or under simulation), as defined in the TMCDB.  In fact, for the Front-End, each Band can load different implementations of driver SW, therefore such changes in harmonic value, or even the tuning algorithm itself, can be achieved by configuring different SW components.  Thus, a heterogeneous mix of antennas with different FE versions will be possible, although operational SW tools (both online and offline) that currently assume that all the antennas have the same hardware model for a given band will no longer be valid. For example, any tools that calculate the expected sensitivity,will need to know the number of upgraded and legacy receiver cartridges in each band along with their tabulated performance levels.

Another long-term issue is how to identify in the raw data (ASDM) which receiver generation (i.e. legacy or upgraded) has been used for an observation.  At least one array-wide mechanism already exists.  The ``receiverId'' field of the ASDM receiver table, which is currently always “0” for all bands, can be incremented to ``1'' for each band once the majority of antennas have been upgraded in that band and are in regular use. However, during the deployment process of a given band, we will need an independent value associated with each antenna.

%%%%%%%%%%%%%%%%%%%
\subsubsection{OSF Data Buffer} % (fold)
\label{osf_data_buffer}

The current OSF Data Buffer is described in Section 5.1 of~\cite{dtas} and Section 3.2.2 of~\cite{CSA-IA}. Fig.~\ref{fig:WSU-OSF-Storage-Hierarchy} shows the data storage and access hierarchy at the OSF.

%%%%%%%%%%%%%%%%%%%
\begin{figure}[htb]
  \centering
    \includegraphics[width=\textwidth]{Figures/WSU-OSF-Storage-Hierarchy.pdf}
  \caption{Data storage and access hierarchy currently at the OSF.}
  \label{fig:WSU-OSF-Storage-Hierarchy}
\end{figure}
%%%%%%%%%%%%%%%%%%%


There are different solutions presented in~\cite{dtas} to improve the performance of the OSF Data Buffer, but they all include:

\begin{itemize}[itemsep=-0.5em]
	\item Upgrade of the base networking to 100~Gbps as default (this includes all network equipment, servers, and storage, if using network-based protocols, such as NFS or iSCSI).
	\item Upgrade of the fiber channel protocols (NGAS backend to storage) to 32~Gbps as default.
	\item Objects are written as a ``file'' to NGAS on a ``per file'' basis.
	\item NGAS GITHUB SW is used, supported NGAS file transfer mechanisms are used (mirroring/subscription).
	\item Discussions are needed to set a disk space availability and data retention policy in order to size the infrastructure.
	\item The assumption that no appreciable impact is foreseen on observation meta-data (Oracle database).
	\item NGAS will provide some load-balancing capability.
	\item Monitoring of new hardware is done by current staff.
	\item The existing QA0+ functionality (running on P-REDOSF) will be replaced by using the calibration pipeline for QA0 (Section~\ref{sec:QA01}).
 %(P-REDOSF) requires minimal refactoring.
%	\item Apart from the QA0+ process, no processing is done at the OSF.
\end{itemize}

The solutions to be down-selected are:

\begin{itemize}[itemsep=-0.5em]
	\item \textbf{Solution 1:} All storage for both NGAS Queue Disk (GAS03), NGAS frontend (FEnn), and NGAS backend are based on Non-Volatile Memory Express (NVME) storage. They all can manage at all times the peak data rate. This is called Scenario~1 in Section 5.3 of~\cite{dtas}.
    %%%%%%%%
	\item \textbf{Solution 2:} Similarly to Solution 1, storage for GAS03 would be based on NVME storage, but for NGAS frontend (FEnn), and NGAS backend storage would be based on slower, cheaper storage. This solution allows for ingesting at the peak data rate until the NVME storage of the GAS03 is filled in, while managing the average data rate at all times. Processing steps that access NGAS@OSF (such as P-REDOSF) only access the NGAS backend storage. This is called Scenario~2 in Section 5.3 of~\cite{dtas}.
    %%%%%%%%%
	\item \textbf{Solution 3:} Variation of Solution 2, where the storage for GAS03 is still based on NVME storage, and there are no more NGAS frontend servers, and instead GAS03 connects directly to NGAS backend. This solution also allows for ingesting at the peak data rate, while managing the average data rate at all times. Processing steps that access NGAS@OSF keep accessing the NGAS backend storage. This makes all data ingestion processes similar, not just ingestion of data from the telescope. The reliability of the NGAS backend servers becomes therefore critical to the online system. This is called Scenario~3 in Section 5.3 of~\cite{dtas}.
    %%%%%%%%%%
	\item \textbf{Solution 4:} Variation of Solution 2, where the storage for GAS03 is still based on NVME storage, while NGAS frontend and backend storage is based on Serially Attached Storage (SAS). This allows for data ingestion process to support the Early WSU peak data rate until the NVME storage of the GAS03 is filled in, while managing the average data rate at all times. The amount of storage is not enough to host the entire archive at the OSF, but in the order of 4 months. This is called Scenario~5 in Section 5.3 of~\cite{dtas}.
    %%%%%%%
	\item \textbf{Solution 5:} Evolution of Solution 4 where the target peak and average data rate are the Later WSU ones. This is called Scenario~6 in Section 5.3 of~\cite{dtas}.
\end{itemize}

Please note that Scenario~4 in~\cite{dtas} is not considered, as it is a mix-and-match which is not clearly defined.

Overall, all solutions support both the Early and Later WSU data rates from the correlator. Performance and reliability modelling including multiple operational use cases will be required to indicate which are the breaking decision points to decide between them beyond their cost, but presently Solutions 4 and 5 seem to provide the right balance between cost and performance, and it is easy to evolve the existing systems in order to provide both solutions.

% subsubsection osf_data_buffer (end)

\subsubsection{OSF to SCO Data Transmission and Buffering} % (fold)
\label{ssub:osf_to_sco_data_transmission_and_buffering}

The current OSF to SCO data transmission is described in~\cite{dtas} Section 6.1 and~\cite{CSA-IA} Section 3.2.4.

There are different solutions given in~\cite{dtas} to improve the 
performance of the OSF to SCO Data Transmission and Buffering to avoid bottlenecks caused by the increased WSU data rates, but they all include:

%%%%%%%%%%%%%%%%%%%
\begin{itemize}[itemsep=-0.5em]
	\item Upgrade the base networking to 100~Gbps as default (this includes all network equipment, servers, and storage, if using network-based protocols, such as NFS or iSCSI).
	\item Other network equipment, such as firewalls, are also upgraded to at least 100~Gbps.
	\item $4 \times$~10~Gbps network links are put in place between OSF and SCO.
	\item No appreciable impact is foreseen on observation meta-data (Oracle database replication).
	\item For the expected data rates, a network backup link (MPLS) is not feasible for peak throughput, so this would impact the storage buffer at the OSF in case of network outage. The current MPLS link (200~Mbps) is kept just for supporting basic services/communications for the Observatory.
\end{itemize}
%%%%%%%%%%%%%%%%%%%

Two solutions are proposed, both only taking into account the physical network route between OSF and SCO:

%%%%%%%%%%%%%%%%%%%
\begin{itemize}[itemsep=-0.5em]
	\item \textbf{Solution 1:} In this solution no backup link is available for science data (meta-data and files). The already existing MPLS network link will only be available for essential Observatory services (telephony, access to applications from either site -- JIRA, Confluence -- and internet service).
	\begin{itemize}[itemsep=-0.5em]
		\item Upgrade the link between OSF-SCO to $4 \times$~10~Gbps. This is above the $\sim 3$~Gbps required by the average WSU data rate.
		\item Proportionally increase storage capacity and speed and the number of NGAS servers.
	\end{itemize}
%%%%%%%%%%%%%%%%%%%

It is expected that by 2025 the $4 \times$~10~Gbps network link is available between OSF and SCO.
	
Not having a back-up link for data transfer of science data implies that the OSF Data Buffer would need to be sized to allow buffer time for network outage contingencies.
	
	\item \textbf{Solution 2:} This solution is a variation of Solution 1, but a parallel backup link is available for science data transfer (both meta-data and files). An additional MPLS network link will be available for essential Observatory services (telephony, access to systems from either site (JIRA, Confluence) and internet service).

	An alternate link provides a different route between OSF-SCO with $4 \times$~10~Gbps capacity.
%	\end{itemize}
	
	As above, this solution expects that by 2025 the $4 \times$~10~Gbps network link is available between OSF and SCO.
	
	Because of the need to provide an alternate route, this solution would require double the operational cost of the first one.
	
	As there will be a back-up link for data transfer of science data, the OSF Data Buffer sizing could be reduced, as the chances of network outage contingencies would decrease.
\end{itemize}

There are no technical risks in any of the solutions, just OPEX and CAPEX costs involved.

% subsubsection osf_to_sco_data_transmission_and_buffering (end)

%%%%%%%%%%%%%%%%%
\subsubsection{Data Transmission from SCO to ARCs and Storage at All Sites} % (fold)
\label{ssub:data_transmission_from_sco_to_arcs_and_storage_at_all_sites}

The current SCO to ARCs data transmission hardware is described in~\cite{dtas} Sections 7.1 and 8.1, and in~\cite{CSA-IA} Section 3.2.5.

The main solution coming from~\cite{dtas} with respect to hardware has to do with moving towards Just a Bunch Of Disks (JBOD) systems, which can currently pack up to 1.7~PB of data, plus error correction overhead up to 2.2~PB, in a 4U unit, which would allow for $\sim 15-17$~PB of data per rack. Using $\sim 9$~racks per site, that would provide on the order of 130~PB, which is assumed to be sufficient for Archive storage and processing. It is also assumed that electrical power of cooling capacity are not limiting factors. Input from ICT, ISOpT and IST is required to validate those assumptions.

In case the available space to place storage and processing equipment not suffice for an ARC, or other constraints like electrical power or cooling capacity be a limiting factor, the current storage system SW could be expanded to allow for a federated deployment at several local sites while at the same time keeping datasets physically collocated. Such a setup would allow an ARC to expand the storage as needed. Care would have to be taken to assure that services that rely on data being collocated would perform satisfactorily. This includes cut-outs (products), CARTA (products), data processing (raw-data), as well as potential science platform access (raw data and products).

Compared to today, we expect a $10\times$ maximum increase of the metadata in the Oracle database. Furthermore, we do not consider that metadata will be a bottleneck. However, larger, faster storage might be required.

%Regarding the network links between SCO and the ARCs, there is currently enough bandwidth (3~Gbps) to transfer the $2\times$BW WSU average data rate, when averaged per cycle. As prices for international bandwidth drop about 16\%/year, double the bandwidth might be feasible at the current operational cost.
Upgrades to the software for the data transfer from SCO to the ARCs, including parallel data transfer for individual files as well as northern-hemisphere data-synchronization between the ARCs through the `partner-site' mechanism of NGAS, will be required.  With these software changes, the bandwidth available today (3~Gbit/s) is already sufficient to transfer the 2$\times$BW WSU average-data rate of estimated 2.1~Gbit/s averaged per cycle. Prices for international bandwidth drop fast at about 16\%/year, which is a factor of two every five years. It can be therefore expected that, even when planning for a bandwidth buffer of a factor of two, the bandwidth cost out of Chile in 2029 will be essentially the same as in 2024.  Extrapolating further, the same cost in 2034 might support 4$\times$BW.

% subsubsection data_transmission_from_sco_to_arcs_and_storage_at_all_sites (end)


\subsubsection{Data Processing Hardware Infrastructure} % (fold)
\label{ssub:data_processing_hardware_infrastructure}

The current Data Processing hardware infrastructure is described in Section~3.3.2 of~\cite{dtas} and in Section~6.1 of~\cite{dpda}.

\paragraph{Storage} % (fold)
\label{dp_hw_storage}

The changes to the data storage infrastructure used by the data processing hardware are discussed in~Section \ref{ssub:data_transmission_from_sco_to_arcs_and_storage_at_all_sites}.

% paragraph dp_hw_storage (end)

\paragraph{Oracle Data Base} % (fold)
\label{par:oracle_data_base}

Based on the relative size between data and metadata in the WSU era it is not expected that the Oracle-based metadata will increase more than ten-fold (Section~8.3 of~\cite{dtas}), and metadata is not expected to be a bottleneck. However, larger, faster storage will be required for the database nodes.

% paragraph oracle_data_base (end)

\paragraph{Processing Cluster} % (fold)
\label{par:processing_cluster}

The current JAOPOST processing cluster will need to be changed to accommodate the processing of the much larger WSU datasets using the new ALMA DPS\index{ALMA Data Processing System (ALMA DPS)} system whose requirements are described in Section~4 of~\cite{dp}. Kepley et al. analyzed the required system performance in~\cite{dpSizeOfCompute}, using the estimated data properties from \cite{drru_mean} (see Section~\ref{subsec:datarates_peak} for an overview). 


%%%%%%%%%%%%%%%%%%%%%%%%
%\begin{table}
%\centering
%\includegraphics[scale=1.1]{Figures/Size-Of-Compute-Table1.pdf}
%\caption{Overview of System Performance Related Quantities for WSU. Copied from Table 2 of~\cite{dpSizeOfCompute}, which assumed 47 operating 12-m antennas. For 50 12-m antennas, the maximum data rate is 1.98~GB/s (see Table~2 of~\cite{drru_peak} for Band~1).}
%\label{computeSizeTable1}
%\end{table}
%%%%%%%%%%%%%%%%%%%%%%%%
%
% TO DO - CHECK/UPDATE BLC NUMBERS
%   - MOVE TO CORRECT PLACE IN APPENDIX, RESTRUCTURE TEXT
%   - DO THE SAME STEPS FOR THE VIS + PROD VOL TABLE.
%   - ADD SYSTEM PERFORMANCE TABLE TO APPENDIX AND DISCUSSION OF SoC
%

\begin{table}[hbt!]
\centering
\begin{talltblr}[%
  caption={Estimated MOUS-level interferometric data rate and required compute system performance.},
  label={tab:computeSizeTable1},
  remark{Note}={\small From \cite{dpSizeOfCompute} and \cite{drru_mean}, with WSU updated to conform with Milestone 5. Note that there is a typo in the units given in the middle 3 rows of Table 2 of \cite{dpSizeOfCompute}, which should be MVis/hr. FLOP/s estimates assume 4000 hr/year of science observing and include the $3\times$ overhead/contingency factor,  as discussed in \cite{computeCostTsi}.}
]{
  colspec={Q[l,m]Q[l,m]*{4}{X[c,m]}},
  width=\textwidth,
  hlines,
  vlines,
  vline{1-2} = {1-2}{0pt},
  hline{1}            = {1-2}{0pt},
  hline{4,5,7,8,10,11} = {2-6}{0pt},
  cells = {font=\small},
  cell{1-2}{3-4} = {my2x, font=\small\bfseries},
  cell{1-2}{5-6} = {my4x, font=\small\bfseries},
}
\SetCell[c=2,r=2]{c}      & & \SetCell[c=2]{c}{WSU at Milestone 5} & & \SetCell[c=2]{c}{pre‑WSU} & \\
                         & & 12 m      & 7 m      & 12 m      & 7 m      \\
\SetCell[r=3]{bg=lightgray}{Data Rate}
  & {Median (GB/s)}                   & $0.136$   & $0.002$   & $5.6\times10^{-3}$   & $2.4\times 10^{-4}$   \\
  & {Time Weighted\\Average (GB/s)}   & $0.477$   & $0.010$   & $9.0 \times 10^{-3}$   & $3.1\times10^{-4}$   \\
  & Maximum (GB/s)                    & $3.48$    & $0.05$    & $4.5\times 10^{-2}$   & $8.7\times 10^{-4}$   \\
\SetCell[r=3]{bg=lightgray}{Visibility Rate}
  & {Median (Gvis/hr)}                & $230$     & $3.2$       & $9.6$   & $0.4$     \\
  & {Time Weighted\\Average (Gvis/hr)}& $822$     & $15$      & $16$   & $0.4$     \\
  & Maximum (Gvis/hr)                 & $6005$    & $77$      & $77$  & $1.3$    \\
\SetCell[r=3]{bg=lightgray}{System\\Performance}
  & {Median (PFLOP/s)}                & $0.07$   & $1.4 \times 10^{-3}$   & $3.0 \times 10^{-3}$   & $1.5\times 10^{-4}$   \\
  & {Time Weighted\\Average (PFLOP/s)}& $0.50$   & $0.021$   & $7.6\times 10^{-3}$   & $5.1\times 10^{-4}$   \\
  & Maximum (PFLOP/s)                 & $10.3$   & $0.13 $   & $9.9 \times 10^{-2}$   & $2.2\times 10^{-3}$   \\
\end{talltblr}
\end{table}



% BELOW IS ALL COMMENTED OUT




\iffalse
\begin{center}
\begin{talltblr}[%
  caption={Overview of system performance-related quantities for WSU.},
  label={tab:computeSizeTable1},
  remark{Note}={Copied from Table 2 of~\cite{dpSizeOfCompute}, which assumed 47 as a realistic number of operating 12-m antennas. For 50$\times$12-m antennas, the maximum data rate is 1.98 GB/s for the Early WSU (see Table~2 of~\cite{drru_peak} for Band 1).}%
]{
  colspec={Q[l,m]Q[l,m] *{6}{X[c,m]}},
  width=\textwidth,
  hlines,
  vlines,
  vline{1-2} = {1-2}{0pt},
  hline{1}   = {1-2}{0pt},
  hline{4,5,7,8,10,11} = {2-8}{0pt},
  cells = {font=\small},
  cell{1-2}{3-5} = {my2x, font=\small\bfseries},
  cell{1-2}{6-8} = {my4x, font=\small\bfseries},
}
\SetCell[c=2,r=2]{c} & & \SetCell[c=3]{c}{WSU at Milestone 5} & & & \SetCell[c=3]{c}{pre-WSU} & &  \\
& & 12 m & 7 m & both & 12 m & 7 m & both \\
\SetCell[r=3]{bg=lightgray}{Data Rate}
  & {Median (GB/s)}                   & $0.136$   & $0.002$   & $0.138$   & $0.060$   & $0.001$   & $0.015$   \\
  & {Time Weighted\\Average (GB/s)}  & $0.477$   & $0.010$   & $0.487$   & $0.220$   & $0.005$   & $0.225$   \\
  & Maximum (GB/s)                    & $3.48$   & $0.05$   & $3.48$   & $1.741$   & $0.026$   & $1.741$   \\
\SetCell[r=3]{bg=lightgray}{Visibility Rate}
  & {Median (Gvis/hr)}                & $230$   & $3$     & $233$    & $103.6$   & $1.6$     & $25.3$    \\
  & {Time Weighted\\Average (Gvis/hr)}& $822$   & $15$    & $837$   & $380.2$   & $7.4$     & $219.0$   \\
  & Maximum (Gvis/hr)                 & $6005$  & $77$    & $6005$  & $3002.8$  & $38.5$    & $3002.8$  \\
\SetCell[r=3]{bg=lightgray}{System\\Performance}
  & {Median (P\ac{FLOP}/s)}                & $0.051$   & $0.001$   & $0.052$   & $0.025$   & $0.001$   & $0.007$   \\
  & {Time Weighted\\Average (P\ac{FLOP}/s)}& $0.364$   & $0.015$   & $0.379$   & $0.183$   & $0.007$   & $0.19$   \\
  & Maximum (P\ac{FLOP}/s)                 & $7.480$   & $0.096$   & $7.480$   & $3.740$   & $0.048$   & $3.740$   \\
\end{talltblr}
\end{center}
\fi 

%\begin{center}
%\begin{talltblr}[caption={Overview of system performance-related quantities for WSU.}, label={tab:computeSizeTable1}, remark{Note}={Copied from Table 2 of~\cite{dpSizeOfCompute}, which assumed 47 as a realistic number of operating 12-m antennas. For 50$\times$12-m antennas, the maximum data rate is 1.98~GB/s for the Early WSU (see Table~2 of~\cite{drru_peak} for Band~1).}]{
%  colspec={Q[l,m]Q[l,m] *{6}{X[c,m]}},
%  width=\textwidth,
%  hlines,
%  vlines,
%  vline{1-2} = {1-2}{0pt},
%  hline{1} = {1-2}{0pt},
%  hline{4,5,7,8,10,11} = {2-8}{0pt},
%  cells = {font=\small},
%  cell{1-2}{3-5} = {my2x, font=\small\bfseries},
%  cell{1-2}{6-8} = {my4x, font=\small\bfseries},
%  }
%\SetCell[c=2,r=2]{c} & & \SetCell[c=3]{c} Early WSU  & & & \SetCell[c=3]{c} Later WSU & &  \\
%& & 12 m & 7 m & both & 12 m & 7 m & both \\
%\SetCell[r=3]{bg=lightgray}{Data Rate} & {Median (GB/s)}& $  0.060$& $  0.001$& $  0.015$& $  0.136$& $  0.002$& $  0.035$\\ 
% & {Time Weighted\\ Average (GB/s)} & $  0.220$& $  0.005$& $  0.127$& $  0.513$& $  0.011$& $  0.296$\\ 
% & Maximum (GB/s)& $ 1.741$& $  0.026$& $ 1.741$& $3.481$& $  0.052$& $3.481$\\ 
%\SetCell[r=3]{bg=lightgray}{Visibility Rate} & {Median (Gvis/hr)}&   103.6&   1.6 &  25.3 & 235.3 & 3.3 & 53.4\\ 
% & { Time Weighted\\ Average (Gvis/hr)} & $  380.2$& $  7.4$& $  219.0$& $ 885.4$& $  15.9$& $  509.4$\\ 
% & Maximum (Gvis/hr)& $3002.8$& $  38.5$& $3002.8$& $6005.5$& $  76.9$& $6005.5$\\ 
%\SetCell[r=3]{bg=lightgray}{System\\Performance} & {Median (P\ac{FLOP}/s)}&   0.025&   0.001 &  0.007 & 0.051 & 0.001 & 0.017\\ 
% & { Time Weighted\\ Average (P\ac{FLOP}/s)} & $  0.183$& $  0.007$& $  0.107$& $ 0.399$& $  0.015$& $  0.233$\\ 
% & Maximum (P\ac{FLOP}/s)& $3.740$& $  0.048$& $3.740$& $7.480$& $  0.096$& $7.480$\\ 
%\end{talltblr}  
%\end{center}

Table~\ref{tab:computeSizeTable1} gives the median, time-weighted average, and maximum values for the data rates, visibility rates, and required system performances calculated using the database of estimated WSU data properties, including estimates for Bands 1 and 2, assembled by Kepley et al. \cite{dpSizeOfCompute}. The weights for the time-weighted average are determined by taking the observing time for each MOUS and dividing by the total observing time for all MOUSs in their sample. These values have been calculated for each of the 50 realizations of their database and then taking the mean of the relevant statistic as the final average result. For comparison, Table~\ref{tab:computeSizeTable2} presents the same values for the current ALMA correlators (BLC or ACA).

%%%%%%%%%%%%%%%%%%%
%\begin{table}
%%x\centering
%\includegraphics[scale=1]{Figures/Size-Of-Compute-Table2.pdf}
%\caption{Overview of System Performance Related Quantities for BLC/ACA. Copied from Table~2 of~\cite{dpSizeOfCompute}.}
%\label{computeSizeTable2}
%\end{table}
%%%%%%%%%%%%%%%%%%%
\begin{center}
\begin{talltblr}[caption={Overview of system performance related quantities for BLC/ACAC.}, label={tab:computeSizeTable2}, remark{Note}={Copied from Table~2 of~\cite{dpSizeOfCompute}. Note the factor of 1000 difference in units compared to Table~\ref{tab:computeSizeTable1}.}]{
  colspec={Q[l,m]Q[l,m] *{3}{X[c,m]}},
  % width=\textwidth,
  hlines,
  vlines,
  vline{1-2} = {1-2}{0pt},
  hline{1} = {1-2}{0pt},
  hline{4,5,7,8,10,11} = {2-5}{0pt},
  cells = {font=\small},
  cell{1-2}{3-5} = {headercolor, font=\small\bfseries},
  }
\SetCell[c=2,r=2]{c} & & \SetCell[c=3]{c} BLC/ACAC & &  \\
& & 12 m & 7 m & both\\
\SetCell[r=3]{bg=lightgray}{Data Rate} & {Median (MB/s)}& $  5.73$& $  0.25$& $  2.29$\\ 
 & {Time Weighted\\Average (MB/s)} & $  9.22$& $  0.32$& $  5.37$\\ 
 & Maximum (MB/s)& $ 45.84$& $  0.89$& $ 45.84$\\ 
\SetCell[r=3]{bg=lightgray}{Visibility Rate} & {Median (Tvis/hr)}& $  9883.0$& $  370.0$& $  3953.0$\\ 
 & {Time Weighted\\Average (Tvis/hr)} & $  15901.0$& $  470.0$& $  9228.0$\\ 
 & Maximum (Tvis/hr)& $ 79067.0$& $  1303.0$& $ 79067.0$\\ 
 \SetCell[r=3]{bg=lightgray}{System\\ Performance} & {Median (T\ac{FLOP}/s)}& $  2.25$& $  0.11$& $  0.84$\\ 
 & {Time Weighted\\Average (T\ac{FLOP}/s)} & $  5.68$& $  0.38$& $  3.39$\\ 
 & Maximum (T\ac{FLOP}/s)& $ 73.86$& $  1.62$& $ 73.86$\\ 
\end{talltblr}  
\end{center}

The general trend of a narrow peak with a long tail, first described in~\cite{drru_mean}, also holds for the system performance calculations. For example, the time-weighted average value for the 12-m~Array is 0.183~PFLOP/s, with a maximum of 3.740~PFLOP/s for Early WSU and 0.399~PFLOP/s with a maximum of 7.480~PFLOP/s for Later WSU. 

Fig.~\ref{fig:Figures_ALMA_WSU_Size_of_Compute_Estimate-Figure_1} shows this trend graphically by showing the probability that the required system performance for a random observation will be larger than the value shown on the x-axis.
%a  observation X has a system performance higher than x. 
%for a given system performance value, $P (X \geq x)$. 
%which is also known as the complementary cumulative distribution. 
From this plot, we see that in Early WSU, approximately 18\% of projects will require a larger system performance than the currently largest BLC/ACA project. This number will increase to 20\% for Later WSU. The time-weighted average required system performance will increase by a factor of 30 (Early WSU) to 70 (Late WSU) above the time-weighted average required system performance for the BLC/ACA.

These calculations are performed with the underlying assumptions below:

\begin{enumerate}
    \item That gridding of the \textit{uv}-data dominates the computational requirements. However, the deconvolution process can dominate the imaging runtimes for ALMA data today, particularly for the largest cubes with emission~\cite{naascMemo121}.
    
    \item I/O considerations are not taken into account. The rates derived in~\cite{dpSizeOfCompute} require that the processing is CPU-bound, not IO-bound. %In other words, the rate at which data is transferred has to be greater than or equal to the rate that it is processed.
    In other words, the total time for disk access - from the storage system as well as the data reading and writing of data to local storage during processing itself - has to be small compared to the entire processing time.
    
    \item Imaging dominates the computational requirements over those of calibration. Calibration of the WSU visibility data is still not a trivial problem. Currently the calibration portion of the ALMA Pipeline is largely serial, thus limited to the performance provided by a single core. Calibrating WSU data will require more efficient processing of calibration data and should include the parallelization of the calibration Pipeline on the EB and likely spw axis.
\end{enumerate}

These estimates are intended as order-of-magnitude estimates. Further modelling will be required for the proper sizing of the infrastructure. Moore's law will make it less costly, and more space efficient to acquire the required systems closer to the date that they are needed, and that other policy decisions might mitigate the size of compute required for the processing cluster.

% paragraph processing_cluster (end)

% subsubsection data_processing_hardware_infrastructure (end)

%%%%%%%%%%%%%%%%%%%%%%%%
\begin{figure}[htb]
  \centering
    \includegraphics[width=\textwidth]{Figures/ALMA WSU Size of Compute Estimate-Figure 1.jpg}
  \caption{The complementary cumulative distribution of the required system performance estimated for the current BLC and the Early and Later stages of the WSU. The y-axis shows the fraction of times that the required system performance will be larger than the value shown on the x-axis. The shaded regions indicate range of the estimates including Bands 1 and 2 with the median indicated as a dotted line. Dotted horizontal lines indicate where the probability of a larger required system performance is 1\%, 5\%, and 10\%. Approximately 18\% of projects will require a larger system performance than the currently largest BLC/ACA project for Early WSU; this value will increase to 20\% for Later WSU.}
  \label{fig:Figures_ALMA_WSU_Size_of_Compute_Estimate-Figure_1}
\end{figure}
%%%%%%%%%%%%%%%%%%%%%%%%

\clearpage


\newpage

%%%%%%%%%%%%%%%%%%%%%%%%%%%%%%%%%%%%%%%%%%%%%%%%%
\subsection{The WSU Software Infrastructure}
%%%%%%%%%%%%%%%%%%%%%%%%%%%%%%%%%%%%%%%%%%%%%%%%%

This section describes how the various software components will evolve, and how new ones will need to be developed, to cope with the data obtained with the WSU. These components are depicted in the right panel of Fig.~\ref{fig:proc_hw_sw_wsu}. A description of the current ALMA software infrastructure is provided in Section~3.3 of~\cite{CSA-IA}.  

%%%%%%%%%%%%%%%%%%%%%%%%%%%%%%%%%%%%%%%%%%%%%%%%%
\subsubsection{Proposals Submission and Project Preparation}
%%%%%%%%%%%%%%%%%%%%%%%%%%%%%%%%%%%%%%%%%%%%%%%%%
\paragraph{The Observing Tool (OT)}\label{subsec:OT}
The PIs submit proposals using the Observing Tool (OT)\index{Observing Tool (OT)} (see Section 3.1.1 of~\cite{CSA-IA} for details). With the WSU, the OT will support setting up the increased number of spws and the available spectral configurations, as well as the calculation of the on-source and total integration times. It will provide a new user interface or SW front-end to make the spectral selection, visualization, and overplot of spectral lines more accessible. The OT will also support the SB generation in taking into account the new observing strategies, including spws, local oscillator definition, correlator modes and correlator efficiency, calibration strategies, and APDM variables (see Section 4.3.1.1 of~\cite{CSA-IA} for details).

\paragraph{Project Tracker (ProTrack)}
The status of any Project, and the associated OUS and SB, is tracked and changed by Project Tracker, or ProTrack\index{ProTrack}. ProTrack will need to adapt to any potential changes to the workflow and status, e.g. per-EB calibration and/or GOUS imaging.
Is also will support automatized work in Phase 2 generation, such as the capability to change the status of multiple SBs at once.

%%%%%%%%%%%%%%%%%%%%%%%%%%%%%%%%%%%%%%%%%%%%%%%%%
\subsubsection{Scheduling and Data Acquisition}
%%%%%%%%%%%%%%%%%%%%%%%%%%%%%%%%%%%%%%%%%%%%%%%%%

\paragraph{Dynamic Scheduling Algorithm (DSA)}
The Dynamic Scheduling Algorithm (DSA) (see Section 3.1.2 of~\cite{CSA-IA}) will continue to serve as the main SW to gather current weather conditions and available resources of arrays and to provide the prioritized list of SBs to be executed and send them for executions. 
With the WSU, resource tracking is a new capability needed for the DSA in order to make maximal use of the ATAC processing resources and allowed peak or mean data rates, particularly when attempting to schedule imaging correlation in one array and VLBI/Beamforming in another array.\footnote{These two modes can be operated simultaneously in separate arrays, but they must share the total available bandwidth in 200~MHz increments, for example 8~GHz each.}
The existence of any Manual Mode Arrays, along with the correlator resources that they are using, also needs to be known to the resource tracker.  
The DSA may also need to be modified to given additional priority boosts to GOUSs with any EBs observed and a higher priority boost to GOUSs with fully observed MOUSs.

\paragraph{Monitoring Tools}
\label{sec:MonitoringTools}
During an SB execution, different monitoring tools will track the system and weather conditions to determine the allowed frequency bands (see Section 3.1.2 of~\cite{CSA-IA}). Array Runtime Manager (ARM) will be the interface for array control and data acquisition, replacing the current OMC, and is expected to be scalable to the WSU. New FEs, BEs, and DTS alarms will be included into an Integrated Alarm System. The Control Command Language (CCL) will be upgraded within the ALMA Common Software (ACS).
Monitoring GUI tools for ATAC and TPGS will provide real-time visualization of autocorrelations and cross-correlations (power and phase) per subscan of the two antennas in a single baseline, polarization, and spw, for a given array. Quicklook will need to accommodate the online calibration results provided by TelCal\index{TelCal}, which will be updated to cope with the WSU capability and calibration strategy. Additionally, the phase-tracker tool and PWV monitor will be modified to support WSU spws.
% %%%%%%%%%%%%%%%
% \begin{itemize}
% %%%%%%%%%%%%%%
% \setlength\itemsep{-0.5em}
%   \item{\textbf{Array Runtime Manager (ARM)}}: will be the interface for Array control and data acquisition, replacing the current OMC. It is expected to be scalable to the WSU. 
%   \item \textbf{Monitoring GUI tools} for ATAC and TPGS will provide real-time visualization of autocorrelations and cross-correlations (power and phase) per subscan of the two antennas in a single baseline, polarization, and spw, for a given Array. 
%   \item  \textbf{Quicklook} provides real time visualization of TelCal\index{TelCal} produced solutions, including per antenna pointing and focus, delay/phase, and atmospheric calibration (opacity, receiver and system temperatures).
%   \item A single instance of \textbf{phase-tracker} is run continuously, and provides running-mean of the phase rms per scan and antenna across EBs.  
%   \item \textbf{Plot\_pwv} provides weather information, phase stability and amount of Precipitable Water Vapor (PWV)\index{Precitable Water Vapor, PWV} in the atmosphere. The latter is taken from all available 12-m antennas and a median value averaged over 10 minutes is displayed. Plot\_pwv also reads the ``dry'' RMS after WVR correction per array produced by phase-tracker and displays these as a time sequence of the past few hours.
% %%%%%%%%%%%%%
% \end{itemize} 
% %%%%%%%%%%%%%
\paragraph{Shiftlog Tool (SLT)}
Shiftlog Tool (SLT) will continue to support logging all observing activities and problems on the 12-m and 7-m Arrays with ATAC and the TP Array with TPGS.


%%%%%%%%%%%%%%%%%%%%%%%%%%%%%%%%%%%%%%%%%%%%%%%%%
\subsubsection{ASDM and APDM}\label{subsec:ASDM-APDM}
%%%%%%%%%%%%%%%%%%%%%%%%%%%%%%%%%%%%%%%%%%%%%%%%%
\paragraph{ASDM}
\label{ASDMchanges}
The ALMA Science Data Model (ASDM) defines the collection of information recorded during an observation that is needed for scientific analysis. Each ASDM contains both binary data files (BDFs) and metadata tables. The metadata tables contain links to other tables and references to the binary files in the Archive. Details about the ASDM can be found in Section~5.8 and 12.2 of~\cite{c10_thb}. 

The ASDM metadata model is quite rich and the underlying implementation already fulfills some important requirements towards supporting the WSU.
However, several additional requirements and changes have been identified for the metadata model to fully support WSU~\cite{dm}:
%%%%%%%%%%%%%%%%%%%%%%%
\begin{itemize}[itemsep=-0.5em]
\item Minor changes in enumerations\footnote{The enumerations is a special type of fixed data, used by the ASDM and APDM, to describe specific elements in the ALMA data model, such as bands, receivers, etc.}, the spw table, and the Correlator Mode table to support both ATAC and TPGS.
%%%%%
\item Additional information may be required from the SB and APDM\index{ALMA Project Data Model} in new tables beyond the SBSummary table to facilitate the future pipeline heuristics.
%%%%%
\item As an alternative, the SB\index{Scheduling Block(s)} and APDM can be made available to the pipeline by archiving them independently from the ASDM, with fetching/export mechanisms similar to ASDM\-export.
\end{itemize}
%%%%%%%%%%%%%%%%%%%%%

The ASDM binary data model is also quite flexible.
However, the high data rate required for the WSU would result in large Binary Large Objects (BLOBs)\footnote{The concept of BLOB  refers to splitting  a single-size binary format dataset into smaller blocks, to facilitate the data streaming. In ALMA, a single BLOB represents the output of the correlators or total-power processor integrations over time, i.e. subcans, spectral-window within a subscan, etc. More details can be found in~\cite{Pokorny2009}.} and may cause problems downstream.
Therefore it may be needed further split the binary BLOB per spw (in addition to the current separation by sub-scan and spectral flow).
An initial evaluation of the required changes in the ASDM binary code indicates that spw separation should be already supported by the current code. With this ``spw separation'' concept in mind, the following requirements/changes for the binary data model are derived:
%%%%%%%%%%%%%%%%%
\begin{itemize}
\setlength\itemsep{-0.5em}
\item The BLOB shall be separated per spw, in addition to per sub-scan and spectral flow.
\item Each BLOB shall be transmitted with a separate bulk data Data Distribution Service (DDS) flow which the archive will store in separate files.
\item The Archive ASDMexport tool shall allow importing a specific set of ASDM BLOB corresponding to a selection of scans, sub-scans and spws.
\end{itemize}
%%%%%%%%%%%%%%%%
   

%%%%%%%%%%%%%%%%
\paragraph{APDM}
The information of any Observing Project is stored in the
ALMA Archive following the ALMA Project Data Model (APDM). More information about the APDM can be found in the Glossary section of the present document and in Section 2.6 of~\cite{CSA-IA}. The current APDM has grown over the years with the basic principle of never renaming or removing any elements so as not to break anything. The downside of this is that some elements of the model are unused, some have misleading names, and others would be misleading if retained for the WSU. An initial APDM draft design for the WSU discussed a few months ago was intentionally not backwards compatible, by a desire for the new APDM to be a clean design and that this would reduce the probability of unforced errors in the new SW.

Irrespective of whether the new APDM is compatible with the old or not, it is unavoidable that SW dealing directly with any of the correlator configuration in the APDM will need to be adapted. What is being discussed is the incidental impact on ALMA SW and database components by the potential introduction of a new APDM that is not compatible with the old model.
Some of this incidental impact, especially on the database side, can be managed with medium effort. The main stumbling block, and where the impact could be very expensive and cause inconvenience to users, is with all SW currently relying on the APDM Castor binding classes~\cite{apdm}. This appears to be most of ALMA SW.

To manage both the current and new (WSU) APDM models using the binding classes, three options have been identified so far:
%%%%%%%%%%%%%
\begin{enumerate}[itemsep=-0.5em]
\item Add the new WSU correlator configuration to the existing model in as clean a way as possible without changing/removing old parts of the APDM.
\item Break with the old APDM model, do a clean WSU design, but use new package names in APDMEntities.
\item Break with the old APDM model, do a clean WSU design, and keep the same package names as in the current APDM.
\end{enumerate}
%%%%%%%%%%%%%


\subsubsection{Control and Data Archiving at the OSF} % (fold)
\label{sec:WSUcontrolAndArchive}

Section~3.3.3 in~\cite{CSA-IA} describes the current ALMA system control, and Section~4.3.4 in the same document describes the WSU impacts. From the impact assessment, no major changes are requried to the overall control and data transfer and archiving architecure up to the OSF. The main salient points to consider are:

\begin{itemize}[itemsep=-0.5em]
	\item \textbf{Bulk data transfer:} not affected by the Early WSU stage, but the component will not be able to meet the requirements at Later WSU stage (with peak rates up to 3.95 GB/s or $\sim 32$~Gbps).
	\item \textbf{Archive client:} requires an upgrade to the current storage capacity to meet the WSU requirements.
	\item \textbf{CORBA RPC:} this technology is used to sahre calibration table slices, and requires huge amounts of memory to handle calibration tables, so more than 64~GiB of RAM are required for each machine involved in the communication.

\end{itemize}

% subsubsection sec:WSUcontrolAndArchive (end)


\subsubsection{Calibration Services}

The telescope calibration services are described in \cite{acsos}. These services are of two types: 
\begin{itemize}[itemsep=-0.5em]
    \item Those performed by observatory staff to measure, update, and maintain all of the  databases necessary for successful science observations and to provide access to those values needed for offline processing, and
    \item Those performed by SW subsystems during and after
observations. 
\end{itemize}

The following subsections describe how ALMA will conduct its calibrations during WSU. 


%%%%%%%
\paragraph{Antenna Calibration}  
%%%%%%%
The WSU system will require all of the same routine measurements
described in Section~3.1.3.1 of \cite{CSA-IA}.  The scripts or
SW packages that are currently used to process them will
need to be adapted to the WSU data format and shape.

%%%%%%%
\paragraph{Calibration Databases}
%%%%%%%
The Jy/K database, the TP OFF position database, and
the source catalog will continue to be used in the WSU
as now (Section~3.1.3.2 of \cite{CSA-IA}).
The wider
bandwidths of the WSU may require populating a new catalog of LO spurious signals and
utilize it for data calibration particularly for TP Array data that does not benefit from the
suppression provided by LO offsetting and 180~degree Walsh switching.

%%%%%%%
\paragraph{Online Calibration}  % relies on receiver
\label{sec:WSUonlineCalibration}
%%%%%%%
Normalization of cross-correlation visibilities will be performed
as currently described in Section~3.1.3.2 of \cite{CSA-IA}.  The
quantization correction for the digitizer will differ in the WSU, as the number of digital bits increases from 3 to 6 (Section~\ref{digitizerSection}).  But the corrections will still be performed in the correlator and total power spectrometer subsystems, using the techniques described in \cite{QC}, including corrections for the various requantization stages in ATAC.  The WVR correction will be applied online only for VLBI/Beamforming mode, for which ATAC
will determine and apply the phasing solutions
rather than the ALMA Phasing System.

Among the calibration tables calculated by TelCal\index{TelCal}, the \Tsys\/ and bandpass calculations are the most impacted.
Although the fundamental technique of \Tsys\/ calibration will remain the same, the cross-correlation data of the \Tsys\/ scans will no longer be stored (to save disk space and speed up data import) and
the spectral resolution of the online calculation by TelCal\index{TelCal} might be coarser than
the science data in some cases due to limited online processing resources or parallelization.
In this case, the full resolution can be regenerated offline, to the extent
that the SNR of the measurement allows it \cite{normalizationAndTsys}.  This regeneration will be done
either by an offline instance of TelCal\index{TelCal}, or by the Pipeline using TelCal\index{TelCal} libraries (TBD). 

Similarly, the spectral resolution of the bandpass calibration
performed by TelCal\index{TelCal} will need to adapt to its current limitations and time allotment.  However, the resolution of these solutions have less 
critical requirements as they are used only by online tools like AosCheck, rather than the Pipeline.

%%%%%%%
\paragraph{Visibility Data Calibration}  
%%%%%%%
Conceptually, all of the same post-observation calibrations currently performed in the interferometric Pipeline \cite{ALMApipelineHeuristics}, as described in Section~3.1.3.3 of \cite{CSA-IA}, will continue to be performed by the DPS.  The main difference is that data will no longer be calibrated on a per-MOUS basis but rather on a per-EB basis (or per-session basis for polarization sessions) \cite{dp} in order to enable the proposed new QA0 process for interferometric data (Section~\ref{sec:QA01}). Other differences in heuristics are highlighted by calibration topic 
in Section~\ref{sec:calibrations}.

%%%%%%%
\paragraph{Total Power Data Calibration}
%%%%%%%
Conceptually, all of the same post-observation calibrations currently performed in the single dish Pipeline\index{Pipeline!Single-Dish} described in Section~3.1.3.3 of \cite{CSA-IA} will continue to be performed by the DPS.  In contrast to the interferometric pipeline, the single dish calibration+imaging Pipeline will likely continue to operate on a per-MOUS basis as it is already has many stages parallelized by EB \cite{dp}.  Further parallelization by spw will be enabled by the proposed changes to the ASDM organization (Section~\ref{ASDMchanges}).

\begin{comment}
\begin{itemize}[itemsep=-0.5em]
\item {\bf Bandpass:} The  spectral index
of calibrators may vary somewhat across the wider WSU spws, which must be accounted for both in the
calibrator source catalog (SC) and in the calibrator
model used for calibration in the DPS. In addition, a recent ancillary study at the NAASC \cite{bandpassMemo}
found that if data volume is a limitation for high spectral resolution science observations, then it
appears safe to observe the bandpass calibrator at spectral resolutions no finer than 1–2~MHz. Additional pre-averaging
in the pipeline bandpass solver in cases of low SNR will likely be necessary as it is now in many cases.

\item {\bf Flux:} The DPS shall calibrate data so that the resulting flux calibration is linear with respect to the flux scaling, and be able to re-scale the flux for individual EBs prior to imaging as necessary. This will allows the per-session calibration to be run as close as possible to the time the data
was acquired, while the fluxes can re-scaled before imaging, if necessary, either because the
source catalog changed or the flux needs to be equalized between sessions.

\item {\bf Gain:} The phase calibrator scans might more often be
observed at coarser spectral resolution (i.e. bandwidth switching) than the science target scans in
order to reduce total data rate.  The implementation of
band-to-band calibration in the pipeline for the 2024 release
includes much of the infrastructure needed to support this
observing mode.

\end{itemize}
\end{comment}
%%%%%%%%%%%%%%%%%%%%%%%%%%%%%%%%%%%%%%%%%%%%%%%%%
\subsubsection{QA0: AosCheck and AQUA/QA0}
%%%%%%%%%%%%%%%%%%%%%%%%%%%%%%%%%%%%%%%%%%%%%%%%%
The quality assurance process is carried out through two main SW tools: the ALMA QUality Assurance (AQUA) and AosCheck. AQUA runs in two different instances, one for QA0 and the other for QA2. Details on these tools can be found in Section 3.3.5 in \cite{CSA-IA}.
%AosCheck is a script running on the results of the calibration performed online by TelCal\index{TelCal}. It produces output (text and plots) displayed by AQUA, which in turn accesses the metadata of the ASDM, to report such parameters as \Tsys\/
%T$_{\rm sys}$ 
%per antenna and scan. AosCheck is mainly a tool to help locate issues at both the antenna level and correlator level and report these to the AoD through AQUA. All t
AosCheck will continue to be used to 
evaluate such parameters as phase stability, system and receiver temperatures, bandpass solutions, and shadowing.  These results are used to
triage EBs that are not worth being calibrated so that they can be promptly rescheduled for re-observation. 
%This first evaluation will be performed only on eight channel-averaged ``channels'' compared to the single channel-averaged data used currently.
%Both AosCheck and 
AQUA QA0 will continue to provide the interface
to QA0 results, which will migrate to being produced by the calibration pipeline (Section~\ref{sec:QA01})

%need re-factoring to take as input an amount of not more than eight channels.

%%%%%%%%%%%%%%%%%%%%%%%%%%%%%%%%%%%%%%%%%%%%%%%%%
\subsubsection{A new Data Processing System (DPS)}
%%%%%%%%%%%%%%%%%%%%%%%%%%%%%%%%%%%%%%%%%%%%%%%%%

The primary recommendation of \cite{dpda} was to design, build, and deliver a new ALMA Data Processing System (ALMA DPS)\index{ALMA Data Processing System (ALMA DPS)} that will meet the requirements of the WSU. Based on the conclusions of \cite{dpda} and the later \cite{dp}, the current ALMA data processing system and data processing operational model are not suitable to handle a significant fraction of the data produced by the WSU, due to the significant increase in the number of channels (and corresponding increased data rates, visibility data volumes, and product sizes). 

The current data processing system already has challenges processing today’s ALMA data. ALMA currently does not image, or images with less than optimal quality, a significant fraction of projects (19\% of the 12-m projects in Cycle 7 \cite{memo623}). A significant increase in the number of channels will only make these challenges more severe. A summary of the properties of the WSU data is provided in Section 2.3.2 of the current document, which is based on \cite{drru_mean} and \cite{drru_peak}.

The current data processing system includes the ALMA Pipeline along with the surrounding infrastructure needed for management of the processing jobs and quality assurance. A list of these and other data processing related sub-systems that will be impacted by the WSU are given in Table 1 of \cite{dp}, including:

%%%%%%%%%%%%%%%%%%%%%
\begin{itemize}[itemsep=-0.5em]
    \item The ALMA Pipeline\index{Pipeline!Interferometry}\index{Pipeline!Single-Dish}, which provides a flexible, intent-driven, framework for processing heterogeneous interferometric and Single Dish ALMA data.
    \item CASA, which provides general purpose tasks for calibrating, imaging, and analyzing radio data.
    \item PLdriver, which executes the pipeline via a queue manager.
    \item ProMo, which manages pipeline products.
    \item CLAI, which launches processing jobs.
    \item AQUA, the primary interface for QA0 and QA2, and 
    \item The Archive, which stores the final data products.
\end{itemize}
%%%%%%%%%%%%%%%%%%%%

A summary of the current ALMA Pipeline and CASA functionality and management structure can be found in Section 3.3.6 of \cite{CSA-IA} and the references therein. Sections 3.3.7.1 through 3.3.7.4 in \cite{CSA-IA} provide a more complete description of the processing infrastructure. See Section 3.1.5 of \cite{CSA-IA} for a description of how the ALMA Pipeline and CASA are used in the QA2 workflow.

From \cite{dp, dpda}, the new ALMA DPS\index{ALMA Data Processing System (ALMA DPS)} needs to:

%%%%%%%%%%%%%%%%%%%%
\begin{itemize}[itemsep=-0.5em]
    \item Make efficient use of the available computing resources.
    \item Enable easier and faster development.
    \item Provide more flexibility and better separation of functionalities (e.g., separate heuristics from data handling), and  
    \item Make greater usage of off-the-shelf software and hardware to support WSU. 
\end{itemize}

Section 4 of \cite{dp} contains a detailed list of ALMA requirements for the design of the new ALMA DPS\index{ALMA Data Processing System (ALMA DPS)}. The key change from the present system is the ability to process smaller chunks of data (e.g., individual scans/spws) in discrete steps (e.g., a single pipeline stage) using the appropriate compute resources for the job rather than requiring all the data in a MOUS to be processed in a monolithic set of steps using a single hardware configuration. This functionality also requires the changes to the ASDM described in Section 3.3.3.1, i.e., splitting of the BLOBs per spectral window (\cite{dp} Requirement \#\,068). The efficient data manipulation tools developed for the new ALMA DPS will also need to be used for interactive data analysis in operations. 

Other required functionality of the new ALMA DPS covers functionality of the existing system, e.g.:

%%%%%%%%%%%%%%%%%%%%%%
\begin{itemize}[itemsep=-0.5em]
    \item A mechanism to initiate and track data processing jobs.
    \item Quality assurance reports.
    \item QA2 interface.
    \item Allocation and management of compute resources, and
    \item Retrieval of data from the Archive, and ingestion into the Archive of data produced by the new ALMA DPS.
\end{itemize}
%%%%%%%%%%%%%%%%%%%%%%

The Data Management and Software (DMS) Division of NRAO is developing a new data processing system for use by both ALMA and ngVLA. The initial ALMA requirements listed in Section 4 of \cite{dp} have been provided to the DMS design team. However, the specific separation of which functionality will be implemented by NRAO DMS and which by ALMA subsystems requires a top down system design and flow down to the necessary requirement. This division of responsibility will thus be an exercise to be agreed between NRAO DMS and ALMA at a later stage.

There will also need to be changes to the ALMA data processing operational model (specifically, QA workflow) to enable the necessary processing efficiency, which are described in \cite{dp}. The key changes, from Section 3.6 of \cite{dp} are:

%%%%%%%%%%%%%%%%%%%%%%
\begin{itemize}[itemsep=-0.5em]
    \item Separation of calibration and imaging;
    \item Calibrating data per-EB (or per-session if polarization);
    \item Maximizing automation of flagging and quality assurance,
    \item Doing appropriate quality assurance at the calibration, MOUS imaging, and GOUS imaging stages;
    \item Producing intermediate, potentially non-science grade, MOUS products;
    \item Final science product imaging and QA at the GOUS (group imaging level) stage.
\end{itemize}

Additional information is also provided in Section~3.1.5 of the present document. 
%%%%%%%%%%%%%%%%%%%%%%

%  ----- OUTDATED TEXT ABOUT DPS BEGINS HERE
% TEXT HAS BEEN REPLACED BY TEXT SENT BY C. VLAHAKIS ON MARCH 27, 2024

%The current ALMA Pipeline provides a flexible, intent driven, framework for processing heterogeneous interferometric (IF) and Single Dish (SD) ALMA data. It currently relies on the underlying CASA SW package, to provide general purpose tasks for calibrating, imaging, and analyzing radio data. A summary of the current ALMA Pipeline and CASA functionality and management structure can be found in Section~3.3.6 of~\cite{CSA-IA} and the references therein. A brief description of how the ALMA Pipeline and CASA are associated to the Quality Assurance Level-2 (QA2) process is provided in Section~3.1.5 of~\cite{CSA-IA} and the references therein. 

%Based on the conclusions of \cite{dpda} and \cite{dp}, the current data processing and data-flow systems are not suitable to handle a significant fraction of the data produced by WSU due to the significant increase in the number of channels. The current data processing system already has challenges processing today's ALMA data. ALMA currently does not image, or images with less than optimal quality, a significant fraction of projects: 19\% of the 12-m projects in Cycle 7 \cite{memo623}. A significant increase in the number of channels will only make these challenges more severe. A summary of the properties of the WSU data is provided in Section~\ref{subsec:datarates_prop} of the current document, which is based on \cite{drru_mean} and \cite{drru_peak}.

%The primary recommendation of \cite{dpda} was to design, build, and deliver a new Data Processing System (DPS) that will meet the requirements of WSU. In particular, from \cite{dp}, the DPS needs to:
%%%%%%%%%%%%%%%
%\begin{itemize}[itemsep=-0.5em]
%    \item Make efficient use of the available computing resources;
%    \item Enable easier and faster development;
%    \item Provide more flexibility and better separation of functionalities (e.g. separate heuristics from data handling); and
%    \item Make greater usage of off-the-shelf SW and hardware to support WSU.
%\end{itemize}

%Section 4 of \cite{dp} contains a detailed list of %ALMA requirements for the design of the DPS. Additional requirements are listed in Sections 4.2.3, 4.2.4, 4.5.3, 6.3, and 6.4 of \cite{uid}. 

%The key change from the present system is the ability to process smaller chunks of data (e.g., individual scans/spws) in discrete steps (e.g., a single pipeline stage) using the appropriate compute resources for the job rather than requiring all the data in a MOUS to be processed in a monolithic set of steps using a single hardware configuration. 

%This DPS software stack is considered to include not just the pipeline heuristics and underlying computational libraries, but also all the surrounding functionality for data processing, namely:
%%%%%%%%%%%%%%%%
%\begin{itemize}[itemsep=-0.5em]
%    \item The mechanism to initiate and track data processing jobs; 
%    \item Quality assurance reports;
%    \item QA2 interface% (currently AQUA);
%    \item Cluster resource management; and
%    \item Retrieval of data from the Archive, and ingestion into the Archive of data produced by the DPS.
%\end{itemize}
%%%%%%%%%%%%%%%%
%The efficient data manipulation tools developed for DPS will also need to be used for interactive data analysis in operations. The division of what functionality will be implemented by NRAO's DMS, and what functionality will be implemented by other ALMA subsystems, requires a top down system design and flow down the necessary requirements.

%There will also need to be changes to the QA2 workflow to enable the necessary processing efficiency, which are described in \cite{dp}. The key changes are:
%\begin{itemize}[itemsep=-0.5em]
%\item Separation of calibration and imaging;
%\item Calibrating data per-EB (or per-session if polarization);
%\item Maximizing automation of quality assurance,
%\item Doing appropriate quality assurance at the %calibration, MOUS imaging, and GOUS imaging stages;
%\item Producing intermediate, non-science grade, MOUS imaging products;
%\item Final science product imaging and QA at the GOUS (group imaging level) stage.
%\end{itemize}

%DPS shall be a drop-in replacement of the current CASA/PL/CLAI functionality, potentially including some of the current TBD functionality of PLDriver. In the interest of preventing any conflict of interest between the DPS functionality developed for ngVLA and what is in the best interest of ALMA, it is vital that the software guiding the data processing workflow, quality assurance, storage and delivery to the user (i.e. at least part of the PL Driver functionality, the state system, ProTrack, AQUA, PROMO and the archive components) remain external to DPS. This will allow development and deployment to be focussed on the needs of ALMA and provide flexibly in light of the rapidly changing requirements leading up to and during the WSU era.

%  ----- OUTDATED TEXT ABOUT DPS ENDS HERE



%%%%%%%%%%%%%%%
\paragraph{Manual Data Reduction}
%%%%%%%%%%%%%%%

Almost all ALMA data will need to be processed with the new Pipeline and any manual interaction, particularly for the processing and delivery of user data, will need to be excluded or significantly limited to corner cases. Current manually processed modes include full-Stokes imaging, VLBI / Beamforming, and Solar. With the WSU, it will be possible to manually process and deliver to PIs some manually processed data, but the data rates will be capped at a more stringent level than for Pipeline processable modes. It is anticipated that, as currently, the primary use case for manual processing will consist of Pipeline scripts with custom tailored arguments but it will also be possible to use non-Pipeline scripts e.g. for Solar or VLBI. For the Solar and VLBI cases, ALMA relies heavily on external expertise to maintain these observing modes, so advance planning will be required and continuous maintenance of the capability could still be at risk if there are extrinsic demands on those resources. This process should occur in coordination with re-commissioning of the observing modes.


%%%%%%%%%%%%
\paragraph{Archival Data}
%%%%%%%%%%%%

It should be possible for any user with typical internet service to access and use the calibrated visibilities for any publicly available, Pipeline-calibrated ALMA dataset from Cycle 3 (TBC) onward, over the life of the ALMA project. It is strongly recommended that the new ALMA DPS be able to automatically reprocess (i.e. both calibrate and image) ALMA Archival data back to Cycle 3 (TBC), assuming that all metadata, flagging, and non-default parameters issues are either fixed in the archival raw data, or captured and provided in an automated way to the ALMA DPS. As a less attractive solution, ALMA could store calibrated visibilities in the Archive after reprocessing with the given current pipeline. The cost tradeoff of the two solutions will be considered at a later stage with the main goal of maximising the scientific value of ALMA's data.

% --------- OUTDATET TEXT ON SUBSECTION CALLED "PREVIOUS DATA". IT HAS BEEN RREPLACED BY TEXT SENT BY C. VLAHAKIS ON MARCH 27, 2024

%It is desirable that the DPS be able to automatically reprocess (i.e. both calibrate and image) ALMA Archival data back to Cycle 3 (TBC), assuming that all metadata, flagging, and non-default parameters issues are either fixed in the archival raw data, or captured and provided in an automated way to the DPS. Alternatively, ALMA could reprocess the earlier data and store the calibrated visibilities in the Archive. The cost tradeoff of the two solutions will have to be considered at a later stage.


% --------- END OF OUTDATET TEXT ON SUBSECTION CALLED "PREVIOUS DATA".  


%%%%%%%%%%%%
\paragraph{Transition Period}
\label{transperiod}
%%%%%%%%%%%%
An initial transition plan was presented in \cite{dpda}. In this tentative plan, once the new imaging pipeline is deployed, the calibration runs would still use the current Pipeline/CASA system, based on MSv2, well until the 2030s. 

A detailed data processing transition plan to go from Pipeline/CASA to deployment of the new ALMA DPS is currently being investigated by a working group who will present results in mid-2024. %Conversion functionality will be required between ASDM and MSv2, as presently, but also between MSv2 and MSv4 as well as ASDM and MSv4 for the DPS components.



%%%%%%%%%%%%%%%%%%%%%%%%%%%%%%%%%%%%%%%%%%%%%%%%%
\subsubsection{AQUA/QA2}
%%%%%%%%%%%%%%%%%%%%%%%%%%%%%%%%%%%%%%%%%%%%%%%%%

AQUA/QA2 is the tool used for performing quality assurance on the ALMA data (both pipeline-able and manually processed). The operational aspects of AQUA/QA2 are described in Section~7.1 of~\cite{dpda}.

Impacts of the WSU on AQUA/QA2 include any associated major changes implemented in the processing life-cycle and workflow, the GOUS imaging, or the QA criteria. Also, AQUA/QA2 will be indirectly impacted by the increment of the data size during WSU, due to its dependence on specific parts of the processing infrastructure and services (i.e. ProMo) that will be affected. These impacts are described in~\cite{dp}.

Specifically, the new WSU high level data processing model will have a significant impact on AQUA/QA2 due to the following:

\begin{itemize}[itemsep=-0.5em]

    \item Per-session calibration

    \begin{itemize}[itemsep=-0.5em]
        \item The QA2 workflow and interface would need some changes to adapt to per-session calibration, but these should be relatively straightforward.
    \end{itemize}

    \item Some form of processing/QA2 assessment at MOUS level with subsequent GOUS imaging and final QA2

    \begin{itemize}[itemsep=-0.5em]
        \item This will completely change the current data processing and QA2 workflow and thus have a major impact on AQUA/QA2, that currently does not deal with groups of MOUSs at all.
        \item Since the technology and design of the current QA2 interface are already quite old anyway it would make sense to design a new AQUA-QA2 interface.
%        \begin{itemize}[itemsep=-0.5em]
%            \item This redesign and reimplementation will be a significant amount of work and a significant part of the FTE we currently have.
%            \item For example, the new QA0 interface (which has the same functionality as the old interface) took $\sim$1 FTE of development, $\sim$6 months of tester FTE, $\sim$6 months subsystem scientist time just to recreate already existing functionality. Redesigning the QA2 part will require at least twice that as the whole workflow and QA2 assessment will change.
%        \end{itemize}
    \end{itemize}

    \item Support automated data processing, quality assurance and ingestion of observatory calibrations.
    \item Adapt ProMo and the staging area to handle and serve significantly larger data products and weblogs.

\end{itemize}

%%%%%%%%%%%%%%%%%%%%%%%%%%%%%%%%%%%%%%%%%%%%%%%%%
\subsubsection{Product Ingestion}
%%%%%%%%%%%%%%%%%%%%%%%%%%%%%%%%%%%%%%%%%%%%%%%%%
\label{sec:ASA}
The software that ingests data-products at SCO, which then get replicated to the ARCs, is called ProductIngestor. This software already can support GOUS-level product ingestions, but will need several updates. First, due to the large product size, larger parallelism in the ingestion process will need to be implemented to increase the efficiency. Second, per-EB data-ingestion will have to be implemented, in coordination with the other relevant subsystems of the ALMA Data Processing Workflow.

%Based on the impact assessment, the only service that would require a redesign is ProductIngestor, as it is only minimally parallelized, able to handle three parallel ingestions at once, which could create a bottleneck. There is no recommendation on the number of parallel ingestions.
%One major change from the current operational model is that ingestions will include GOUS-level science products. However, the current design is expected to allow for it.

%The Archive will need to support EB-level, session-level, MOUS-level and, one major change from the current operational model, GOUS-level ingestion, as well as searching and retrieval.

%%%%%%%%%%%%%%%%%%%%%%%%%%%%%%%%%%%%%%%%%%%%%%%%%
\subsubsection{User Interface with the Data (UID)}
%%%%%%%%%%%%%%%%%%%%%%%%%%%%%%%%%%%%%%%%%%%%%%%%%

% ALMA provides support to the user community through the ALMA Regional Centers (ARCs) and ARC nodes. Support is provided from proposal preparation to data analysis, including access to ALMA data through the ALMA Science Archive. The large data volumes that will come with the WSU will require enhancing the support to the user community.

ALMA provides support to the user community for data access and analysis through the ALMA Regional Centers (ARCs) and ARC nodes. Access to the data is primarily through the ALMA Science Archive. Besides the ``raw'' visibilities, users can download calibration tables, scripts to generate the calibrated visibilities, image cubes, continuum images, and web logs. Available archive services to facilitate archival use include search capabilities by observing parameters, a machine learning based text similarity search for proposal and publication abstracts, and object type search for all SIMBAD and NED observations falling into footprints of ALMA observations. The Science Archive also allows users to query and download data programmatically. The suite includes a cut-out service, allowing users to extract portions of all ALMA FITS images and cubes without having to download the full data.

With the WSU, the typical size of a visibility dataset will increase by $\sim7\times$ for the Early WSU and $\sim18\times$  for the Later WSU~\cite{drru_mean}. In the most extreme cases, the data volume will increase by 690$\times$ over current dataset sizes. Overall, 15-35\% of the visibility data (from Early to Late WSU) will be larger than the largest visibility data volumes obtained with the current system (see Section~\ref{sub:wsu_data_rates}). Therefore, an increasing number of users will need observatory support to handle the large data volumes produced by the WSU.

The ALMA Data Processing System shall provide image analysis capabilities that enable the Observatory to generate (and deliver to users) an overview of the astronomical signals contained in a dataset, and enable users to perform in-depth analyses on selected subsets of data, with both automated and interactive use cases.

Table~\ref{tab:uid} summarizes the existing services that will need to be enhanced and the new services that are needed to optimize the scientific output from the WSU as described in \cite{uid}. Users will retain the ability to reprocess data on their own computers with the WSU,\footnote{Caveat: subject to the processing limitations of the ALMA Data Processing System.} but archive processes can be parallelized and automated to increase the efficiency in handling the increased data volume. As an example, users should be able to request a specific field or individual spw(s) as opposed to downloading all visibility data for a given MOUS. 

The user interface to WSU data can be enhanced by providing new services. Computing infrastructure that can be accessed remotely to recalibrate data, and to analyze images and visibility data, will enable users with modest computing resources to analyze larger datasets. Reimaging services will allow users to submit requests to image datasets with customized parameters to meet their specific science goals. Advanced data products have the potential to greatly facilitate data mining by providing lists of sources in two dimensional images, detected lines in cubes, extracted spectra, moment maps of detected lines, and flux measurements of sources and lines. These new services will expedite the scientific analysis of datasets in the WSU era.


\begin{center}
\begin{talltblr}[caption={Enhanced and new user services to support WSU science.},label=tab:uid,]{
  colspec={QX},
  width=\textwidth,
  hlines,
  vlines,
  cells={font=\linespread{0.5}\small},
  row{1} = {bg=headercolor, font=\small\bfseries},
  rowsep=4pt,
  }
   Service &
   Benefit to the user community\\
\SetCell[c=2]{c,bg=highlight2}{Enhancements to existing services/products}\\
Science Archive &
\begin{minipage}{10cm}
\begin{itemize}[leftmargin=0.3cm,rightmargin=0.3cm,itemsep=0em]
\item Enable parallel downloads of individual files.
% \item Integration into re-calibration, science enabling infrastructure, and re-imaging services
\item Augment ADMIT to handle larger data volumes efficiently.
\item Periodic automated reprocessing of the ALMA holdings by DMG with the latest version of the
   pipeline.
\end{itemize}
\end{minipage}\\
Restoration of calibrated data & 
\begin{minipage}{10cm}
\begin{itemize}[leftmargin=0.3cm,rightmargin=0.3cm,itemsep=0em]
\item Request calibrated measurement sets from Cycle 3 onwards.
\item Request calibrated data for specific source or spectral window.
\item See Sections 7.3 and 7.4 in \cite{uid} for a full set of requirements.
\end{itemize}
\end{minipage}\\
Data Processing and Analysis Software & 
\begin{minipage}{10cm}
\begin{itemize}[leftmargin=0.3cm,rightmargin=0.3cm,itemsep=0em]
\item Provide software that can calibrate, image and combine ALMA data from Cycle 3 onward on user's hardware.
\item Identify spatio-spectral regions of interest.
\item Provide data-analysis functionality (potentially in a separate package).		
\end{itemize}
\end{minipage}\\
CARTA & 
\begin{minipage}{10cm}
\begin{itemize}[leftmargin=0.3cm,rightmargin=0.3cm,itemsep=0em]
\item Maintain current functionality of CARTA, scaled to the data sizes of the WSU.
\end{itemize}
\end{minipage}\\
%
Science Portal & 
\begin{minipage}{10cm}
\begin{itemize}[leftmargin=0.3cm,rightmargin=0.3cm,itemsep=0em]
\item Develop modern interface, complying with state-of-the-art requirements on user-friendliness/usability, accessible from all browsers and devices.
\end{itemize}
\end{minipage}\\
%
%
\SetCell[c=2]{c,bg=highlight2}{New services}\\
Science enabling infrastructure   &
\begin{minipage}{10cm}
\begin{itemize}[leftmargin=0.3cm,rightmargin=0.3cm,itemsep=0em]
\item Offer an interactive and batch analysis collaborative platform that is accessible remotely.
%\item Allow users to re-calibrate and re-image ALMA data, plus perform data combination, image analysis, and visibility analysis
\item Allow users to perform image and visibility analysis.
\item  See Sections 8.3 and 8.4 in \cite{uid} for a full set of requirements.
\end{itemize}
\end{minipage}\\
%
Reimaging service         & 
\begin{minipage}{10cm}
\begin{itemize}[leftmargin=0.3cm, rightmargin=0.3cm,itemsep=0em]
\item Provide interface to users for re-imaging ALMA data.
\item Allow imaging of line+continuum, continuum subtracted, and continuum only.
\item Allow user to select MOUS\index{Member Observing Unit Set}, spectral windows, weighting schemes, and other imaging parameters.
\item See Sections 9.3 and 9.5 in \cite{uid} for a full list of requirements.
\end{itemize}
\end{minipage}\\
%
Advanced data products     & 
\begin{minipage}{10cm}
\begin{itemize}[leftmargin=0.3cm,rightmargin=0.3cm,itemsep=0em]
\item Generate advanced data products to facilitate scientific research.
\item Include source lists detected in 2D images, detected lines in cubes, extracted spectra, moment maps of identified lines, and flux measurements of identified sources and lines.
\item See Section 10.3 in \cite{uid} for a full list of requirements.
\end{itemize}
\end{minipage}\\
\end{talltblr}
\end{center}
